\documentclass[11pt,a4paper]{article}
\usepackage[hmargin=27mm,vmargin=25mm]{geometry}
\usepackage[T1]{fontenc}
\usepackage{lmodern,amsmath,amssymb,amsthm,mathtools,microtype}
\usepackage[dvipsnames]{xcolor}
\usepackage[colorlinks=true,linkcolor=MidnightBlue,citecolor=MidnightBlue,urlcolor=MidnightBlue,breaklinks=true]{hyperref}
\usepackage{bookmark}
\hypersetup{pdftitle={Pointwise Binary Correlations in Abelian Splitting Semigroups},pdfauthor={Roy van Rijn}}
\newtheorem{theorem}{Theorem}[section]
\newtheorem{proposition}[theorem]{Proposition}
\newtheorem{lemma}[theorem]{Lemma}
\newtheorem{corollary}[theorem]{Corollary}
\theoremstyle{definition}
\newtheorem{definition}[theorem]{Definition}
\theoremstyle{remark}
\newtheorem{remark}[theorem]{Remark}
\newcommand{\Z}{\mathbb Z}
\newcommand{\Q}{\mathbb Q}
\newcommand{\C}{\mathbb C}
\newcommand{\N}{\mathbb N}
\newcommand{\E}{\mathbb E}
\newcommand{\Prob}{\mathbb P}
\newcommand{\1}{\mathbf 1}
\DeclareMathOperator{\supp}{supp}
\DeclareMathOperator{\lcm}{lcm}
\DeclareMathOperator{\rad}{rad}
\DeclareMathOperator{\Res}{Res}
\DeclareMathOperator{\Cl}{Cl}
\DeclareMathOperator{\Gal}{Gal}
\DeclareMathOperator{\Frob}{Frob}
\DeclareMathOperator{\TV}{TV}
\newcommand{\ee}[1]{e^{2\pi i #1}}
\allowdisplaybreaks[2]
\setlength{\emergencystretch}{1.5em}
\title{Pointwise Binary Correlations in Abelian Splitting Semigroups}
\author{Roy van Rijn}
\date{}
\begin{document}
\maketitle
\begin{abstract}
We study representations of an integer as a sum of two positive integers
whose prime factors satisfy fixed abelian splitting restrictions.
For semigroups defined by subgroups of a fixed reduced residue group,
we establish pointwise asymptotics with explicit arithmetic factors and
a positive logarithmic saving, uniformly in the target. The resulting
congruence criterion characterizes eventual representability.
In particular, integers supported on primes congruent to one modulo
twelve have an explicit binary representation asymptotic of order
\(N/(\log N)^{3/2}\) on every sufficiently large admissible target.
The proof adapts Walsh's multiplicative-invariance framework, using
the recent zero-free and weighted-dilation theorems of OpenAI,
a density model retaining the relevant congruence characters, and
transference for reflected correlations with growing targets.
We count numerical values without representation multiplicities.
As an arithmetic application, we derive explicit leading asymptotics
for sums of primitive norm values in every fixed imaginary quadratic order.
\end{abstract}
\section{Introduction}
How many positive integers \(n\le X\) have the property that both
\(n\) and \(n+1\) are sums of two squares? Write \(b_d(n)\) for
the indicator of the positive values of \(x^2+dy^2\), with
\(x,y\in\Z\), and put \(b_d(n)=0\) for \(n\le0\).
Coordinates are unrestricted and may vanish; each numerical value
is counted once. Under the analytic inputs stated below, we obtain
\begin{equation}\label{eq:ck-headline}
 \#\{1\le n\le X:b_1(n)=b_1(n+1)=1\}
                         \sim\frac{X}{2\log X}.
\end{equation}
This is the shift-one Connors--Keating prediction. It belongs to
Smilansky's five-form family of pair-correlation predictions.

\begin{corollary}\label{cor:smilansky-five-pairs}
Under the inputs of Section~\ref{sec:inputs}, for
\(d\in\{1,2,3,4,7\}\) and fixed \(h\ne0\),
\[
 \sum_{1\le n\le X}b_d(n)b_d(n+h)
       =\mathscr T_{d,h}\frac X{\log X}+o_{d,h}(X/\log X),
\]
with precisely the coefficients predicted in
\cite[Section 6.3, equations (6.2)--(6.4)]{smilansky}.
In particular,
\begin{equation}\label{eq:minus-eight-shift-one}
 \sum_{n\le X}b_2(n)b_2(n+1)
                       \sim\frac1{\sqrt2}\frac X{\log X}.
\end{equation}
\end{corollary}

The coefficients and their exceptional-prime conditions are evaluated
in Section~\ref{sec:quadratic-pairs}. In particular these are ordinary
value counts, with no weights for the number of coordinate representations.

\subsection{A second headline: genus governs primitive value counts}

Fix one negative fundamental discriminant \(D\) and two fixed
discriminants \(\Delta_i=f_i^2D\). Let \(Q_i,Q'_i\) be primitive
positive definite integral binary forms of discriminant \(\Delta_i\),
in the same genus for each \(i=1,2\). Thus both slots belong to
orders in the one quadratic field \(\Q(\sqrt D)\); no mixed-field
assertion is made. If \(P_Q\) is the set of positive values of
\(Q\) at coprime coordinates and
\(R_{Q_1,Q_2}(N)=\#\{(a,b)\in P_{Q_1}\times P_{Q_2}:a+b=N\}\),
then the classical route gives
\begin{equation}\label{eq:genus-headline}
 R_{Q_1,Q_2}(N)=R_{Q'_1,Q'_2}(N)
                  \left(1+O((\log N)^{-1/5})\right)
\end{equation}
on sufficiently large common locally admissible targets.
The two slots may have different genera and conductors; replacement
is within each respective discriminant and genus. Under the recent
analytic inputs we also obtain an explicit mixed leading term of
order \(NJ_S(N)/\log N\), with relative error \(o(1)\).
There is no class-number divisor: a numerical value can be represented
by every class in its compatible genus. Section~\ref{sec:forms}
states the general results before giving the ideal-class and local-filter
proofs; Section~\ref{sec:norms} recovers the principal specialization.

The counting conventions and strengths are different:
\begin{center}\small
\begin{tabular}{@{}>{\raggedright\arraybackslash}p{.19\linewidth}>{\raggedright\arraybackslash}p{.34\linewidth}>{\raggedright\arraybackslash}p{.41\linewidth}@{}}
\hline
Result family&What is counted&Conclusion\\\hline
CK/Smilansky&Ordinary values; unrestricted coordinates&Fixed-shift asymptotic with the stated little-\(o\) error\\[1mm]
Subgroup/coset sums&Integers with prescribed prime support&Target-uniform asymptotic with a logarithmic saving\\[1mm]
Primitive-form sums&Values admitting coprime coordinates in two fixed orders of one field&Explicit leading term with relative \(o(1)\)\\[1mm]
Genus replacement&Comparison of two primitive-form value counts&Relative logarithmic saving through the classical route\\\hline
\end{tabular}
\end{center}
The first three asymptotic families use the imports in
Section~\ref{sec:inputs}; genus replacement does not.
A \emph{primitive form} \([a,b,c]\) has \(\gcd(a,b,c)=1\).
A \emph{primitive representation} \(Q(x,y)\) has
\(\gcd(x,y)=1\). \emph{Coprime summands} in \(a+b=N\)
means \(\gcd(a,b)=1\), a separate condition imposed only when stated.
For example, \(4=2^2+0^2\) is an ordinary sum of two squares,
but any representation of four has both coordinates even and is
not primitive. Consequently the primitive-form theorem does not
extend the ordinary CK statement to arbitrary forms. Its class-error
saving also does not supply a logarithmic remainder for the full
conductor-prefix leading term.

\subsection{The common mechanism, in the CK example}

The central result is the indicator-to-model comparison of
Theorem~\ref{thm:model-comparison}. Multiplication by suitable primes
preserves an indicator and its arithmetic model. One-variable means
center their difference, and a positive sparse sieve controls exceptional
prime-factor records. A shared graph reduction then leaves two frequency
estimates: one for fixed shifts and one for reflected targets. Their
assembly replaces the indicators in the binary count by the models.
The model correlation can then be evaluated from local arithmetic.

Here is the model for \(b_1\). At outer scale \(X\), put
\(L=\log X\), \(y=\exp(L^{2/5})\),
\(\mathcal P_y=\{p:p\le y\}\), and
\(V_y=\prod_{p\le y}(1-1/p)\). Write every \(n\) uniquely as
\(n=rv\), where \(r=n_{\mathcal P_y}\) contains its full small-prime
powers and \(v\) has no prime factor at most \(y\).
The rough ordinary-value series is
\[
 F_y(s)=\prod_{\substack{p>y\\p\equiv1\ (4)}}(1-p^{-s})^{-1}
         \prod_{\substack{p>y\\p\equiv3\ (4)}}(1-p^{-2s})^{-1}
       =(s-1)^{-1/2}H_y(s).
\]
For a sufficiently deep fixed \(K\), its smooth density expansion is
\[
 M(v)=\sum_{j=0}^K\frac{H_y^{(j)}(1)}{j!}
            \frac{(\log v)^{-1/2-j}}{\Gamma(1/2-j)},\qquad
 \mathcal D(rv)=b_1(r)\bigl(1+\chi_{-4}(v)\bigr)\frac{M(v)}{V_y}.
\]
These formulas apply on the polynomial rough scales used in the proof;
the smooth bounded extension is specified in
Section~\ref{sec:model-preparation} and the quadratic application.
The leading term of \(M\) is \(H_y(1)/\sqrt{\pi\log v}\).
Dividing by \(V_y\), the density of rough integers, gives the
conditional density on the rough variable. Every allowed rough value
has \(\chi_{-4}(v)=1\). The correction \(1+\chi_{-4}(v)\)
retains this bias; average density alone would miss a leading contribution.
For a split prime \(p\le y\), multiplication by \(p\) changes only
\(r\), preserves \(b_1(r)\), and leaves both \(v\) and
\(\mathcal D\) unchanged. Replacing
\(b_1(n)b_1(n+h)\) by \(\mathcal D(n)\mathcal D(n+h)\)
with an error \(o(X/\log X)\) is the analytic task. The remaining
model product gives the coupled local factors and the coefficient
\(1/2\) in \eqref{eq:ck-headline}.

\subsection{Semigroups and the reflected-target obstacle}

The same mechanism reaches a different geometry. Let
\(C_{12}=\{n\ge1:p\mid n\Rightarrow p\equiv1\pmod{12}\}\),
including one. Every sufficiently large \(N\equiv2\pmod{12}\)
is a sum of two elements of \(C_{12}\), and
the ordered value-pair count
\(R_{12}(N)=\#\{(a,b)\in C_{12}^2:a+b=N\}\) satisfies
\begin{equation}\label{eq:opening}
 R_{12}(N)=A_{12}J_{12}(N)\frac{N}{(\log N)^{3/2}}
                  \left(1+O((\log N)^{-\delta})\right),
 \qquad\delta>0.
\end{equation}
The positive factors are explicit in Theorem~\ref{thm:c12}.
For \(C^-_4=\{n\ge1:p\mid n\Rightarrow p\equiv3\pmod4\}\),
every sufficiently large even target likewise has an explicit
asymptotic of order \(N/\log N\), with a positive logarithmic saving.
The subgroup and coset theorems allow fixed arithmetic data and retain
all surviving characters, including the complex branches forced by
the prime class \(3\pmod{10}\). Numerical saving exponents are
collected in Section~\ref{sec:assembly}.

One-variable density does not control the particular line \(a+b=N\),
and averages can conceal exceptional targets. The nonconstant linear
parts of \(n\) and \(N-n\) are dependent, so the independent-linear-form
correlation theorems of Loughran--Matthiesen do not directly apply
\cite[Theorems 1.9 and 3.10]{loughranmatthiesen}.
The reflected frequency argument must retain cancellation and
normalization when the target enters the rational denominators.

Iwaniec--Pomyka\l a's quartic-norm work is a direct antecedent
\cite{iwaniecpomykala}. For \(\Q(\zeta_{12})\), its splitting conditions
select the primes \(1\pmod{12}\). Their Theorem~2 gives conditional
pointwise order bounds for coprime sums; hypothesis (P4) is uniform
in growing shifts. Their Theorem~1 treats fixed differences
unconditionally. They count numerical values once.
Our unrestricted leading term gives a coprime asymptotic by M\"obius
inversion, with the extra target factor in
Corollary~\ref{cor:coprime-c12}. Historical priority beyond these
precise comparisons remains separate from the proof.

\subsection{Attribution, dependencies and scope}

Walsh's multiplicative-invariance framework treats positive integer
divisor orders and a fixed nonzero shift \cite{walsh}.
He explicitly proposes sums-of-two-squares indicator correlations
in Section~1.2.4, pp.~5--6. The application idea is already present
there; the work here supplies the complete indicator models,
arithmetic hypothesis checks and local factors, together with the
uniform reflected-target argument. Exact common-prime invariance,
centered one-variable rough means and a positive sparse envelope are
the independently checkable assumptions of the comparison theorem.
No centered binary cancellation estimate is assumed.

The recent imports are the weighted-dilation root-moment theorem
\cite[Section 3]{weightedgraphs} and the all-conductor Dirichlet
zero-free half-plane \(\Re s>11/12\)
\cite[Part I, Theorem 3.1]{zerofree}. The stronger Part II claim is
not needed. Walsh states that classical estimates could recover his
fixed-shift divisor theorem \cite[p.~2]{walsh}, while his printed
Lemma~5.1 and Section~6.4 still use the recent zero-free input.
Our fixed-shift moduli are polynomial-logarithmic, but their Mellin
heights reach \(x(\log x)^C\). Section~\ref{sec:fixed-classical-boundary}
proves a classical intermediate-height estimate and isolates the
remaining high-frequency obligation. Both recent inputs remain in
the fixed-shift theorem. A future replacement there would still leave
the reflected proof's growing-conductor issue and the graph input.

For primitive values in the fixed field \(\Q(\sqrt D)\), the classical
Kloosterman, weighted-correlation and positive-restoration route gives
pointwise order bounds. The class-error sieve then proves
\eqref{eq:genus-headline} independently of the recent inputs.
Its full classical core argument is Appendix~\ref{sec:older-route}.
The common mixed local-filter calculation gives the explicit leading
coefficient under the recent binary theorem; the principal norm case
is a specialization. Constants and thresholds depend on the fixed
forms and orders, and no numerical onset or complete small-exception
classification is asserted here. The exact classifications and
certificates belong to the companion article, and the deposited
window paper addresses where summands may be placed.

\subsection{Notation}

\begin{center}\small
\begin{tabular}{@{}>{\raggedright\arraybackslash}p{.28\linewidth}>{\raggedright\arraybackslash}p{.66\linewidth}@{}}
\hline
Symbol&Role\\\hline
\(X\), \(L=\log X\), \(h\)&Outer scale, its logarithm, and the fixed shift\\
\(\mathcal P_y\), \(n_{\mathcal P_y}\), \(V_y\)&Small primes, complete small-prime part, and rough-integer density\\
\(f_i,\mathcal D_i,\xi_i\)&Indicators, arithmetic models, and centered errors\\
\(D,D_i\), \(U_i\)&Auxiliary prime products and their dyadic lengths\\
\(H_i\), \(\mathscr A\), \(\mathcal E\)&Residue subgroups, a prime coset, and forbidden prime sets\\
\(B_0\), \(S\)&Finite exceptional-prime sets in the ordinary-indicator and primitive-form applications\\
\(\Delta_i=f_i^2D\)&Order discriminants sharing one fixed fundamental discriminant \(D\); here \(f_i\) denotes a conductor\\
\(P_Q,\Lambda_Q\), \(G\)&Primitive numerical values, their local shadow, and the applicable class group\\\hline
\end{tabular}
\end{center}
The modulus and the arithmetic data stay fixed. Definitions in each
application specify its groups and exceptional primes; \(h\) is
reserved for the additive shift throughout.

\section{The correlation theorem and congruence admissibility}
Throughout, \(L_0\ge2\) is even,
\(G=(\Z/L_0\Z)^\times\), and \(H_1,H_2\le G\) are fixed.
Define
\[
 \mathcal P_i=\{p\nmid L_0:p\bmod L_0\in H_i\},\quad
 C_i=\langle\mathcal P_i\rangle,\quad
 \rho_i=\frac{|H_i|}{\varphi(L_0)},\quad
 \kappa=2-\rho_1-\rho_2.
\]
An empty product is one. Fix \(U\subset H_1\), \(V\subset H_2\),
and coprime positive \(L_0\)-supported integers \(A,B\).
Write
\[
 R_{U,V}^{A,B}(N;I)=\#\{Ax+By=N:x\in C_1,\ y\in C_2,
      \ x\bmod L_0\in U,\ y\bmod L_0\in V,\ Ax/N\in I\}.
\]
Omit \(I\) when \(I=[0,1]\). Put
\begin{align}
 c_i&=\frac{1}{\Gamma(\rho_i)}
    \lim_{s\downarrow1}(s-1)^{\rho_i}
              \prod_{p\in\mathcal P_i}(1-p^{-s})^{-1},\label{eq:marginal}\\
 \Gamma_N&=\frac{L_0}{|H_1||H_2|}
    \#\{(u,v)\in U\times V:Au+Bv\equiv N\pmod{L_0}\},\label{eq:finite}\\
 \mathcal E&=\{p\nmid L_0:p\notin\mathcal P_1,
                                  \ p\notin\mathcal P_2\},\notag\\
 C_{\mathcal E}&=\prod_{p\in\mathcal E}
               \left(1-\frac{1}{(p-1)^2}\right),\qquad
 J_{\mathcal E}(N)=\prod_{\substack{p\mid N\\p\in\mathcal E}}
                              \frac{p-1}{p-2}.\label{eq:singular}
\end{align}
The infinite product in \eqref{eq:singular} converges absolutely and
is positive, since \(2\mid L_0\). The constants \(c_i\) are positive.
For an explicit convergent expression, let \(r_i(p)>1\) denote the
order of \(pH_i\) in \(G/H_i\) at a forbidden good prime. Then
\begin{equation}\label{eq:marginal-product}
 c_i=\frac{(\varphi(L_0)/L_0)^{\rho_i}}{\Gamma(\rho_i)}
 \prod_{\substack{\eta\in H_i^\perp\\\eta\ne1}}L(1,\eta)^{\rho_i}
 \prod_{\substack{p\nmid L_0\\p\notin\mathcal P_i}}
                     (1-p^{-r_i(p)})^{1/r_i(p)}.
\end{equation}
Characters here are literal modulo \(L_0\). Continuation along the
positive real axis from \(s>1\) fixes the powers; conjugate factors
are paired. The Euler identity proving this formula is given in
Section~\ref{sec:model}.

\begin{theorem}\label{thm:general}
With all these data fixed, let
\(\delta=1/(8192\varphi(L_0))\). Uniformly in \(N\) and every interval
\(I\subset[0,1]\),
\begin{align*}
 R_{U,V}^{A,B}(N;I)
 &=c_1c_2\Gamma_N C_{\mathcal E}|I|
       \frac{NJ_{\mathcal E}(N)}{AB(\log N)^\kappa}\\
 &\quad+O\left(\frac{NJ_{\mathcal E}(N)}
                         {AB(\log N)^{\kappa+\delta}}\right).
\end{align*}
If \(\Gamma_N=0\), the count is identically zero.
For \(\Gamma_N>0\) and \(I=[0,1]\), the error is a relative
\(O((\log N)^{-\delta})\). More generally it is relative \(o(1)\)
when \(|I|(\log N)^\delta\to\infty\).
\end{theorem}

The implicit constant and threshold may depend on all fixed data.
No numerical threshold, uniformity in a growing modulus, or
power-short interval assertion is included. For noncoprime fixed
coefficients, first divide by their gcd when it divides \(N\).
The even-modulus convention makes the prime two a bad prime throughout
the analytic theorem. Passing from an odd modulus to its even multiple
requires retaining this convention; freely admitting the prime two is
a separate valuation-profile operation, not a change of notation.

\begin{theorem}\label{thm:c12}
Let \(\mathcal E_{12}=\{p\nmid6:p\not\equiv1\pmod{12}\}\), and set
\[
 A_{12}=\frac{4\pi}{\Gamma(1/4)^2}
       \sqrt{\frac{\log(2+\sqrt3)}{6}}
       \prod_{p\in\mathcal E_{12}}
                        \left(1-\frac{2}{p(p-1)}\right),\qquad
 J_{12}(N)=\prod_{\substack{p\mid N\\p\in\mathcal E_{12}}}
                                      \frac{p-1}{p-2}.
\]
For all sufficiently large \(N\equiv2\pmod{12}\),
\[
 R_{12}(N)=A_{12}J_{12}(N)\frac{N}{(\log N)^{3/2}}
                 \left(1+O((\log N)^{-1/32768})\right).
\]
All other residue classes have zero count.
\end{theorem}
\begin{proof}[Deduction from Theorem~\ref{thm:general}]
Take \(L_0=12\), \(H_1=H_2=U=V=\{1\}\), \(A=B=1\).
Then \(\rho_i=1/4\) and \(\Gamma_N=12\) on the indicated class.
The four characters give the exact identity
\begin{align*}
 F_{12}(s)^4={}&\zeta(s)L(s,\chi_{-4})L(s,\chi_{-3})L(s,\chi_{12})\,
        (1-2^{-2s})(1-3^{-2s})
        \prod_{p\in\mathcal E_{12}}(1-p^{-2s})^2.
\end{align*}
Here the three nonprincipal characters displayed are primitive.
Use \(L(1,\chi_{-4})=\pi/4\),
\(L(1,\chi_{-3})=\pi/(3\sqrt3)\), and
\[
 L(1,\chi_{12})=\int_0^1\frac{1-x^4}{1+x^6}\,dx
 =\left[\frac{1}{2\sqrt3}
 \log\frac{x^2+\sqrt3x+1}{x^2-\sqrt3x+1}\right]_0^1
 =\frac{\log(2+\sqrt3)}{\sqrt3}.
\]
Consequently
\[
 c_{12}=\frac{1}{\Gamma(1/4)}
 \left(\frac{\pi^2\log(2+\sqrt3)}{54}
       \prod_{p\in\mathcal E_{12}}(1-p^{-2})^2\right)^{1/4}.
\]
In \(12c_{12}^2C_{\mathcal E_{12}}\), combine the factors by
\((1-p^{-2})(1-(p-1)^{-2})=1-2/(p(p-1))\).
This is exactly \(A_{12}\). The real density integrates to one;
there is no extra beta-integral factor.
\end{proof}

The uniformity in the target also permits a coprimality restriction.
Only the primes allowed in the summands can divide their gcd.
Removing them changes the leading term by the following finite product.
\begin{corollary}\label{cor:coprime-c12}
Put
\[
 R_{\mathrm{cop}}(N)=\#\{(a,b)\in C_{12}^2:a+b=N,\ (a,b)=1\},
 \qquad
 P_{12}(N)=\prod_{\substack{p\mid N\\p\equiv1\ (12)}}(1-p^{-1}).
\]
Uniformly over sufficiently large \(N\equiv2\pmod{12}\),
\[
 R_{\mathrm{cop}}(N)
 =A_{12}J_{12}(N)P_{12}(N)\frac{N}{(\log N)^{3/2}}
       \left(1+O((\log N)^{-1/65536})\right).
\]
Every other residue class has zero count. In particular, every
sufficiently large admissible target has coprime summands.
\end{corollary}
\begin{proof}
M\"obius inversion gives the exact identity
\begin{equation}\label{eq:coprime-inversion}
 R_{\mathrm{cop}}(N)
 =\sum_{\substack{d\mid N\\d\in C_{12}}}\mu(d)R_{12}(N/d).
\end{equation}
Indeed, division by any common divisor preserves both prime supports,
and every common divisor lies in \(C_{12}\). Moreover \(d\equiv1\pmod{12}\),
so \(N/d\equiv2\pmod{12}\), and
\(J_{12}(N/d)=J_{12}(N)\).

Write \(L=\log N\), \(\ell=\log\log(3N)\),
\(P=P_{12}(N)\), and \(\delta=1/32768\).
We use the elementary divisor estimates
\begin{align}
 S_0&=\sum_{\substack{d\mid N\\d\in C_{12}}}\frac{|\mu(d)|}{d}
       =\prod_{\substack{p\mid N\\p\equiv1\ (12)}}(1+p^{-1})
       \ll\ell,\qquad P^{-1}\ll\ell,\label{eq:coprime-products}\\
 S_1&=\sum_{\substack{d\mid N\\d\in C_{12}}}
                          \frac{|\mu(d)|\log d}{d}
       =S_0\sum_{\substack{p\mid N\\p\equiv1\ (12)}}
                                     \frac{\log p}{p+1}
       \ll\ell^3.\label{eq:coprime-logmoment}
\end{align}
For \eqref{eq:coprime-products}, split the prime products at \(L\).
Mertens' product bound controls primes at most \(L\) by \(O(\log L)\);
the remaining factors have bounded product, since
\(\sum_{p\mid N,\ p>L}p^{-1}\le1/\log L\).
For \eqref{eq:coprime-logmoment}, the primes at most \(L\) contribute
at most \(\sum_{n\le L}(\log n)/n\ll(\log L)^2\), and the others
contribute at most \(L^{-1}\sum_{p\mid N}\log p\le1\).

Split \eqref{eq:coprime-inversion} at \(d=\sqrt N\).
The bound \(R_{12}(M)\le M\) shows that its large-divisor part is
\(O(\sqrt N\,\tau(N))\), where \(\tau\) is the divisor function.
For \(d\le\sqrt N\), Theorem~\ref{thm:c12} applies uniformly, and
\[
 (L-\log d)^{-3/2}
       =L^{-3/2}+O(L^{-5/2}\log d).
\]
Consequently, using \(\sum_{d\mid N,\ d\in C_{12}}\mu(d)/d=P\)
and extending the main sum back to all divisors,
\begin{align*}
 R_{\mathrm{cop}}(N)
 ={}&A_{12}J_{12}(N)\frac{N}{L^{3/2}}
       \left(P+O(L^{-\delta}S_0+L^{-1}S_1
                            +N^{-1/2}\tau(N))\right)\\
 &\quad+O(\sqrt N\,\tau(N)).
\end{align*}
After division by the main term, the error is
\[
 O\!\left(L^{-\delta}\ell^2+L^{-1}\ell^4
                  +N^{-1/2}\tau(N)L^{3/2}\ell\right).
\]
Here \(J_{12}(N)\ge1\) and \(\tau(N)=N^{o(1)}\).
The displayed error is \(O(L^{-\delta/2})\), proving the claim.
\end{proof}

\begin{remark}[Comparison with the quartic-norm target weight]
The weight in \cite[Theorem 2]{iwaniecpomykala} is
\(W_{12}(N)=\prod_{p\mid N}(1-\varepsilon_p/p)\), with
\(\varepsilon_p=1\) on allowed primes and \(-1\) otherwise.
On the admissible class,
\[
 \frac{J_{12}(N)P_{12}(N)}{W_{12}(N)}
 =\frac23\prod_{\substack{p\mid N\\p\equiv5,7,11\ (12)}}
                 \frac{p(p-1)}{(p-2)(p+1)}\asymp1
\]
uniformly in \(N\), since the factors are \(1+O(p^{-2})\).
The factor \(2/3\) comes from \(2\mid N\); the prime three does
not divide an admissible target. Thus the coprime asymptotic has
the same order as the earlier conditional coprime bounds.
\end{remark}

We next clarify the congruence criterion without analytic inputs.
Write \(C_H(U)=\{n\in C_H:n\bmod L_0\in U\}\).
\begin{lemma}\label{lem:finite-images}
A residue \(r\bmod M\) is represented by \(C_H(U)\) if and only if
its reduction modulo \(\gcd(M,L_0)\) is that of some \(h\in U\),
and every forbidden prime \(p\mid M\), \(p\nmid L_0\), satisfies
\(p\nmid r\).
\end{lemma}
\begin{proof}
Necessity follows from prime support. Choose such an \(h\), and for
each allowed \(p^k\Vert M\) put \(e_p=\min(v_p(r),k)\).
Set \(D=\prod p^{e_p}\), so \(D\in C_H\) and \((D,L_0)=1\).
We seek a prime \(\ell\equiv hD^{-1}\pmod{L_0}\).
Where \(e_p<k\), impose
\[
 \ell\equiv(r/p^{e_p})(D/p^{e_p})^{-1}\pmod{p^{k-e_p}}.
\]
Where \(e_p=k\), impose \(\ell\equiv1\pmod p\).
At primes of \(M\) not dividing \(D\), impose the unit condition
\(\ell\equiv rD^{-1}\pmod{p^k}\). These congruences are compatible
with the \(L_0\)-condition precisely by the stated hypotheses.
They give a reduced CRT class. Dirichlet's theorem supplies a prime
\(\ell\) outside the finite exceptional set; then \(D\ell\) has
the required residue and state.
\end{proof}

\begin{theorem}\label{thm:local}
The condition \(\Gamma_N>0\) is equivalent to solvability of
\(Ax+By\equiv N\pmod M\) with actual
\(x\in C_{H_1}(U),y\in C_{H_2}(V)\), for every \(M\ge1\).
For all sufficiently large \(N\), it is also equivalent to positive
additive representability.
\end{theorem}
\begin{proof}
Taking \(M=L_0\) proves necessity. For sufficiency choose a root
counted by \(\Gamma_N\) and work modulo \(K=\lcm(M,L_0)\).
At a prime dividing \(L_0\), at least one coefficient is a unit;
adjust its variable to lift the equation while preserving the fixed
unit states. At a prime outside \(L_0\), both coefficients are units
and \(p\ge3\). Of the \(p\) solutions modulo \(p\), requiring a
unit on each forbidden side removes at most two. A surviving pair
lifts to every higher power. CRT and Lemma~\ref{lem:finite-images}
realize the required two residues by actual semigroup values.
Finally Theorem~\ref{thm:general} is positive for all sufficiently
large targets with \(\Gamma_N>0\).
\end{proof}

\begin{remark}[A joint congruence obstruction]\label{rem:coupled}
Take \(L_0=12\), \(H=\{1,11\}=\ker\chi_{12}\).
Then \(H+H=\{0,2,10\}\pmod{12}\), so \(N\equiv4\pmod{12}\)
is never a sum. Nevertheless \(H\) projects onto the full unit groups
modulo four and three. The separate closures are \(\Z_2^\times\)
and \(\Z_3^\times\); at other primes they are \(\Z_p\) or
\(\Z_p^\times\), according as that prime is allowed or forbidden.
Two units sum to every element of \(\Z_p\) for odd \(p\), and to
every even element for \(p=2\). Thus these targets pass every
separate \(p\)-adic test. The missing condition is the simultaneous
residue state modulo twelve.
\end{remark}

\section{Analytic inputs and the comparison framework}\label{sec:inputs}

The analytic argument has three sources with different roles.
Walsh's density-model and multiplicative-invariance framework
\cite[Sections 2--7]{walsh} supplies the architecture of the comparison.
He explicitly anticipates the sums-of-two-squares indicator application
in Section 1.2.4 and states on p.~2 that classical estimates could
replace the zero-free input in his fixed-shift divisor theorem.
Our completed fixed-shift proof below retains that input at large
Mellin heights; Section~\ref{sec:fixed-classical-boundary} separates
the classical range from the remaining estimate.
The zero-free theorem \cite[Theorem 3.1]{zerofree} supplies a fixed
half-plane for Dirichlet \(L\)-functions. The weighted-dilation theorem
\cite[Section 3]{weightedgraphs} transfers an independent-label operator
estimate to a moment on integers. We state the latter input with its
normalizations, since constancy on a mark fiber and the order of parameter
choices matter in this application.

\subsection{The zero-free input}

\begin{proposition}[Imported zero-free theorem]\label{prop:zf}
Every Dirichlet \(L\)-function, including \(\zeta\), is nonzero in
\[
                         \Re s>11/12.
\]
The pole at \(s=1\) of a principal character is allowed. Both conductor
and height are unrestricted.
\end{proposition}

This is the Dirichlet conclusion of
\cite[Theorem 3.1, printed p.~10]{zerofree}. The same theorem treats
finite-order Hecke characters over \(\Q(\sqrt{-3})\); that stronger
form is not needed here. Neither the boundary line nor the sharper
\(7/8\) assertion in Part II of that source is used. The input is
nonvanishing on the open half-plane, not an asserted numerical
threshold for the representation theorem. Estimates for individual
\(L\)-functions used below require the usual functional equation,
Euler product, and holomorphic logarithm on zero-free rectangles in
addition to this proposition.

\subsection{Prime groups and the two operators}

For this subsection only, write
\[
                \Lambda=\log x,\qquad T=\log\Lambda.
\]
Fix constants
\[
 0<\alpha<\beta<\gamma<\delta<0.47,\quad
 c_0,C_0,v_-,v_+>0,\quad 0<\vartheta<1,\quad \ell\ge1.
\]
There are disjoint collections \(\mathcal G_s,\mathcal G_b\) of
pairwise disjoint finite prime sets \(P_g\), with
\[
 c_0T\le|\mathcal G_s|,|\mathcal G_b|\le C_0T.
\]
They satisfy
\[
\begin{split}
 e^{\Lambda^\alpha}\le p\le e^{\Lambda^\beta}
       &\quad(g\in\mathcal G_s,\ p\in P_g),\\
 e^{\Lambda^\gamma}\le p\le e^{\Lambda^\delta}
       &\quad(g\in\mathcal G_b,\ p\in P_g),\\
 v_-\le V_g:=\sum_{p\in P_g}\frac1p\le v_+,
       &\qquad \mu_g(p)=\frac1{V_gp}.
\end{split}                                                    \tag{3.1}\label{eq:inputbanks}
\]
The sets need not be intervals. This freedom is needed when taking
identity residue classes and deleting primes dividing the target.

Fix a sufficiently large integer \(J\), and put \(M=J+\ell\).
A pattern \(\nu\) specifies subsets
\(C_g\subseteq\{1,\ldots,J\}\), with
\[
 C_g=\varnothing\quad(g\in\mathcal G_s),\qquad
 |C_g|\le m\quad(g\in\mathcal G_b),\qquad t_g=\ell+|C_g|.
\]
The slots in \(\{1,\ldots,J\}\setminus C_g\) are shared between the
two endpoint lists. The remaining \(t_g\) slots at each endpoint
are unshared. Form \(D\) from the shared labels and, in the
\(C_g\)-slots, fresh independent draws with laws \(\mu_g\).
Thus \(D\) has \(J\) prime factors from each group, counted with
multiplicity. Write \(D=D_sD_b\) according to group size.

Each pattern has a complex coefficient \(K_\nu\). It may depend on
the ordered big-group source labels, target labels, and auxiliary
labels used in \(D_b\), and on no small-group labels. Assume
\[
                 \sum_\nu\sup |K_\nu|\le\Lambda^{C_1},           \tag{3.2}\label{eq:inputmass}
\]
where \(m,C_1\) are fixed independently of \(J\). All operators below
average source and target slot permutations independently in each
applicable group. These averages are orthogonal projections.

The \emph{ideal operator} acts on
\[
 \mathcal H_{\mathrm{id}}=
 L^2\!\left(\prod_{g\in\mathcal G_b}P_g^M,
                    \bigotimes_{g\in\mathcal G_b}\mu_g^{\otimes M}\right).
\]
Numerical repetitions are permitted. Its row operation retains shared
labels, samples the unshared target labels and auxiliary labels
independently, and multiplies by
\[
                         K_\nu D_b^{\,i\zeta}e(\Theta D_b),
             \qquad e(t)=\exp(2\pi it).
\]
Sum over patterns and compose with the permutation projections on
both sides; denote the result by \(\mathcal T(\Theta)\).

A \emph{physical state} consists of an integer \(n\) and an ordered
\(M\)-tuple in every group, with distinct labels within each group,
all dividing \(n\). Give every such state at \(n\) mass
\(\prod_gV_g^{-M}\), and use counting measure on the integer position.
This defines \(\mathcal H_{\mathrm{ph}}\). Put
\[
 \omega_g(n)=\#\{p\in P_g:p\mid n\},\quad
 \omega(n)=\sum_g\omega_g(n),\quad
 s=|\mathcal G_s|+|\mathcal G_b|.
\]
For a fixed integer \(k\) with \(0<|k|\le\Lambda^{C_2}\), the
\emph{physical operator} \(A\) moves from \(n\) to \(n'=n+kD\).
It sums over the physical target lists with the prescribed shared
labels, using row normalization \(\prod_gV_g^{-t_g}\); auxiliary
labels retain their independent probability laws. Every auxiliary
label is forbidden from both endpoint lists, although auxiliary
labels may equal each other. Multiply the row by
\[
 K_\nu D^{\,i\zeta}
       \vartheta^{(\omega(n)-Ms)/2}
       \vartheta^{(\omega(n')-Ms)/2},                            \tag{3.3}\label{eq:inputphysical}
\]
sum over patterns and apply all endpoint permutation projections.
The adjoint reverses the displacement and conjugates the coefficient
with the endpoint roles interchanged. The joint endpoint
normalization in each group is \(V_g^{-M-t_g}\) in either direction.

\begin{proposition}[Imported absolute bounds]\label{prop:graphabsolute}
Replace every elementary transition coefficient of \(A\) by its
absolute value before summing patterns and averaging permutations.
The resulting operator and its adjoint have row sums at most
\(\Lambda^{C_{\mathrm{abs}}}\). The exponent may depend on \(m,C_1\)
and the fixed group data, but is independent of \(J\).
The same assertion holds with integer position replaced by independent
uniform residues at all group primes.
\end{proposition}

This is \cite[Lemma 3.4]{weightedgraphs}. The independence from \(J\)
is part of the assertion. To see its source, once the shared labels
are fixed, the new target slots in a group cost at most
\[
       \sup_{z\ge0}V_g^{-t_g}(z+t_g)^{t_g}\vartheta^{z/2}.
\]
Here \(t_g\le\ell+m\). This is a fixed constant per group and its
product over \(O(T)\) groups is a fixed power of \(\Lambda\).
Reversing the joint endpoint measure gives the column bound.

\subsection{Root moments, residue replacement, and endpoints}

\begin{proposition}[Imported local transference]\label{prop:rootmoment}
For every fixed \(E>0\), there is \(E_{\mathrm{id}}>0\), depending
only on \(E\) and the fixed group data, such that the following holds.
Suppose
\[
                   \sup_{\Theta\in\mathbb R}
                        \|\mathcal T(\Theta)\|_{2\to2}
                              \le\Lambda^{-E_{\mathrm{id}}}.
                                                                    \tag{3.4}\label{eq:inputideal}
\]
After \(m,C_1\) have been fixed, take a sufficiently large fixed \(J\).
Put \(R=\lceil\Lambda^{0.52}\rceil\). For every fixed \(\eta>0\)
and \(I=[X,2X)\), with \(x^\eta\le X\le x^{1/\eta}\),
\[
 \frac1{|I\cap\Z|}\sum_{n\in I\cap\Z}
       \langle u_n,(AA^*)^Ru_n\rangle
                         \le\Lambda^{-2RE}.                    \tag{3.5}\label{eq:inputroot}
\]
Here \(u_n\) is one on every physical mark list at position \(n\)
and zero elsewhere. The threshold may depend on all the fixed
parameters, including \(J,C_2,\eta\). No bound on \(\zeta\) is
required beyond \eqref{eq:inputideal}.
\end{proposition}

This is \cite[Theorem 3.5]{weightedgraphs}. It imposes return of the
integer position. The final mark list need not equal the initial
list, and \(u_n\) is not normalized. Substituting a trace with
identical initial and final mark lists would change the statement.

The arithmetic replacement inside that theorem has the following
precise form, which also explains its scale limitation.
For a fixed compatible path with \(r=2R\) edges, let \(h_i\) be
the integer offsets and \(\mathcal A_i\) its active prime set at
visit \(i\). Set \(\vartheta_i=\vartheta\) at internal visits and
\(\vartheta_0=\vartheta_r=\sqrt{\vartheta}\).
The position-dependent factor is
\[
 W(n)=\prod_{i=0}^{r}
       \left(\prod_{p\in\mathcal A_i}1_{p\mid n+h_i}\right)
       \vartheta_i^{\,\#\{p\notin\mathcal A_i:
                     \ p\text{ is a group prime},\ p\mid n+h_i\}}.
                                                                    \tag{3.6}\label{eq:inputpath}
\]

\begin{proposition}[Imported residue replacement]\label{prop:graphcrt}
For every fixed \(B>0\), the average of \eqref{eq:inputpath} on
\(I\cap\Z\) equals its expectation with independent uniform residues
at all group primes, multiplied by \(1+O(\exp(-BrT))\).
An incompatible collection of active congruences has zero weight
in both models.
\end{proposition}

This is \cite[Lemma 3.6]{weightedgraphs}. For clarity, its relative
form retains all deterministic damping belonging to active primes.
The product \(Q_{\mathrm{act}}\) of active primes satisfies
\[
              \log Q_{\mathrm{act}}
                   =O_J(rT\Lambda^\delta)=o(\Lambda).
\]
After those factors are extracted, a Bonferroni expansion of degree
\(O_B(rT)\) controls the remaining damping. Its independent
expectation is at least \(\exp(-O(rT))\); its total interval
counting error is \(\exp(-\eta\Lambda+o(\Lambda))\).
The inequality \(0.52+\delta<1\) makes that error negligible
even after division by the active congruence density. Together
with Proposition~\ref{prop:graphabsolute}, this permits summation
over paths with absolute coefficients. These observations describe
the imported replacement and do not assert replacement on intervals
shorter than its stated range.

\begin{corollary}[Imported endpoint consequence]\label{cor:graphendpoint}
Under Proposition~\ref{prop:rootmoment}, let \(f\in\mathcal H_{\mathrm{ph}}\)
be supported on positions \(I=[X,2X)\), constant on the chosen mark
list at each position, with value \(f_n\). If
\[
 \sup_{n\in I}|f_n|\le\exp(O(\sqrt{\Lambda})),\qquad
          \|f\|,\|g\|\le X^{1/2}\Lambda^{C_3},
\]
then
\[
                    |\langle f,Ag\rangle|
                        \ll X\Lambda^{-E+2C_3}.                 \tag{3.7}\label{eq:inputendpoint}
\]
Only \(f\) must be constant on mark lists.
\end{corollary}

This is \cite[Corollary 3.11]{weightedgraphs}. Its proof partitions
\(I\) into intervals on the displacement scale
\(\exp(O_J(T\Lambda^\delta))\), uses spectral H\"older for
\(AA^*\), and bounds the positive matrix
\((\langle u_n,(AA^*)^Ru_m\rangle)\) by its trace.
It makes no equal-fiber-mass assumption.

\subsection{What must be proved for the binary application}

The propositions above do not state a binary correlation theorem.
Walsh's theorem \cite[Theorem 1.1]{walsh} concerns
\(d_k(n)d_l(n+h)\) for positive \emph{integer} divisor orders and a
fixed nonzero shift \(h\). His Section 1.2.4 identifies other
indicator correlations as possible extensions. The membership
indicators and the reflected line in the present article require
separate arguments.

There are three interfaces. First, the one-variable Euler series
has a fractional branch \((s-1)^{-\rho_i}\). Its density model must
retain every character trivial on \(H_i\); an average scalar density
would lose genuine main terms. Second, the independent-label
comparison must be constructed for the fixed progression filters
and for prime groups in an identity residue class. The coefficient
bound \eqref{eq:inputmass} and the restriction to big labels must
hold before the final \(J\) is chosen. In particular, the source's
ideal-kernel construction for unrestricted interval prime sets is
not silently imported for these progression filters. Third, the
relation is \(n+(N-n)=N\), with the target growing on the summation
scale. Its comparison terms, endpoint norms, and localization
errors must be uniform in that target. We prove these interfaces
in the following sections; none is a consequence of the root
moment alone.

The parameter order is consequently part of the proof: choose the
required final accuracy, choose the ideal accuracy supplied by
Proposition~\ref{prop:rootmoment}, construct the signed family and
its fixed \(m,C_1\), and only then choose \(J\). The threshold in
\(x\) is taken last. The modulus and subgroup data remain fixed
throughout.

\section{Subgroup models and application checks}\label{sec:model}
We verify the arithmetic hypotheses of Theorem~\ref{thm:model-comparison}
for the subgroup restrictions, then evaluate their model correlation.
We first work with a single proper subgroup \(H\le G\), of density
\(\rho\), and write \(f=\1_{C_H}\), \(k_H=[G:H]=1/\rho\).
The construction for the two sides is separate. Put
\[
 L=\log N,\quad y=\exp(L^{2/5}),\quad S=\{p:p\le y\},\quad
 E=\exp(\sqrt L),\quad V_S=\prod_{p\le y}(1-1/p).
\]
Let \(n_S\) denote the \(S\)-supported part of \(n\), and
\(n_{S'}=n/n_S\). Define
\[
 E_+(y)=\prod_{\substack{p\le y\\p\in\mathcal P_H}}(1-1/p),\quad
 E_-(y)=\prod_{\substack{p\le y\\p\nmid L_0,\ p\notin\mathcal P_H}}(1-1/p),
 \quad \alpha_0=\varphi(L_0)/L_0.
\]
Then \(V_S=\alpha_0E_+(y)E_-(y)\) for large \(N\).

\subsection{Fractional poles and centered rough means}
Character orthogonality gives, for \(\Re s>1\),
\begin{equation}\label{eq:euler-branches}
 F_H(s)=\prod_{p\in\mathcal P_H}(1-p^{-s})^{-1}
       =G_H(s)\prod_{\eta\in H^\perp}L(s,\eta)^\rho.
\end{equation}
At a forbidden good prime of coset order \(r>1\), the correction
is \((1-p^{-rs})^{1/r}\); thus \(G_H\) converges absolutely for
\(\Re s>1/2\). The literal principal character carries the bad
Euler factors, proving \eqref{eq:marginal-product}. Every
\(\eta\in H^\perp\) equals one on the support of \(f\); every such
twist therefore has the same pole branch.

Define the rough series and its regular part by
\[
 F_R(s)=F_H(s)\prod_{\substack{p\le y\\p\in\mathcal P_H}}(1-p^{-s})
        =(s-1)^{-\rho}H_R(s).
\]
Here \(H_R(1)=\Gamma(\rho)c_HE_+(y)\). Fixed-order Cauchy estimates
on a disk of radius \(c/\log y\) give
\begin{equation}\label{eq:rough-derivatives}
 |H_R^{(j)}(1)|/j!\ll_j H_R(1)(C\log y)^j.
\end{equation}
For a fixed depth \(K\), define on polynomial scales in \(N\)
\begin{equation}\label{eq:hankel-model}
 M_{R,K}(t)=\sum_{j=0}^K\frac{H_R^{(j)}(1)}{j!}
                 \frac{(\log t)^{\rho-j-1}}{\Gamma(\rho-j)}.
\end{equation}
For \(N^a\le t\le N^C\), with fixed \(0<a<C\), this is
\(c_HE_+(y)(\log t)^{\rho-1}(1+O_K(L^{-3/5}))>0\).
Extend it smoothly and nonnegatively below that range with the same
upper bound. Only the polynomial bulk enters principal means.

\begin{lemma}\label{lem:rough-means}
For fixed \(B,A_0>0\), sufficiently large fixed \(K\) has the
following property. If \(b\) is bounded and periodic of period at
most \(L^B\), and \(W\) is a smooth dyadic weight at a polynomial
scale \(T\), with fixed derivatives bounded by fixed powers of \(L\),
then
\begin{align*}
 &\sum_n W(n)b(n)f(n)\1_{(n,\prod_{p\le y}p)=1}\\
 &\quad-\sum_n W(n)b(n)\1_{(n,\prod_{p\le y}p)=1}
       \left(\sum_{\eta\in H^\perp}\eta(n)\right)
                  \frac{M_{R,K}(n)}{V_S}\ll T L^{-A_0}.
\end{align*}
The same holds with a factor \(n^{iu}\), \(|u|\le L^B\).
\end{lemma}
\begin{proof}
Refine the period to \(Q=\lcm(q,L_0)\) and expand on unit residues
by characters. All primes dividing \(Q\) are at most \(y\).
For \(\psi\bmod Q\), the twisted rough series is
\[
 \prod_{\eta\in H^\perp}L(s,\psi\eta)^\rho\,G_\psi(s)
 \prod_{\substack{p\le y\\p\in\mathcal P_H}}(1-\psi(p)p^{-s}).
\]
If \(\psi\) is inflated from \(H^\perp\), this equals \(F_R\);
otherwise all its character factors are nonprincipal.
Use smooth Mellin inversion at height \(\exp(L^{1/10})\), moving
to \(1-c/\log y\). The classical zero-free region is wider there
at these heights and polylogarithmic conductors. Siegel's theorem
places a possible exceptional real zero outside the strip for large
\(N\), after choosing its fixed exponent depending on \(B\).
The finite Euler product costs \(L^{O(1)}\); the displacement gives
\(\exp(-cL^{3/5})\), dominating polynomial \(L\)-function growth at
that height. Smooth tails are smaller than every fixed logarithmic
power. The threshold can be ineffective.

Hankel integration gives \eqref{eq:hankel-model} for each principal
branch. Its relative Taylor error is \(O_K(L^{-3(K+1)/5})\),
by \eqref{eq:rough-derivatives}.
Choose \(K\) after the derivative and character-count costs.
For the unweighted rough population the series
\(\zeta(s)\prod_{p\le y}(1-p^{-s})\) has principal density \(V_S\),
and only the trivial character contributes. The projection
\(\sum_{\eta\in H^\perp}\eta\) therefore reproduces every principal
branch of the first sum with its exact coefficient. The factor
\(n^{iu}\) changes only the derivative budget and choice of \(K\).
\end{proof}
For \(U\subset H\), set
\begin{equation}\label{eq:corrected-model}
 D_U(n)=k_H\1_{n\bmod L_0\in U}f(n_S)
                       \frac{M_{R,K}(n_{S'})}{V_S},\qquad
 \xi_U(n)=f(n)\1_{n\bmod L_0\in U}-D_U(n).
\end{equation}
Then \(D_U\ge0\), \(D_U\ll L^{-3(1-\rho)/5}\) and
\(|\xi_U|\ll1\). Conditioning on an exact record \(n_S=r\le E\)
leaves residual state \(U/r\), so Lemma~\ref{lem:rough-means}
centers its principal mean. Fixed divisor-growth record coefficients
cost only powers of \(L\), since
\(\sum_{r\ S\text{-supported}}d_j(r)/r\ll L^{O_j(1)}\).
Rankin's bound with exponent \(c/\log y\) gives
\begin{equation}\label{eq:record-tail}
\#\{n\le T:n_S>E\}\ll TL^{O(1)}e^{-cL^{1/10}}
 \quad(N^a\le T\le N^C).
\end{equation}
Indeed the Euler moment is \(L^{O(1)}\), whereas
\(E^{-c/\log y}\) gives the exponential saving. The same argument
permits fixed divisor weights. For \(H=G\), use the exact model
\(D_U=\1_{n\bmod L_0\in U}\) and \(\xi_U=0\).

\subsection{The model correlation}
Use separate constructions for the two subgroups with common \(S\).
\begin{proposition}\label{prop:model-correlation}
For a smooth \(0\le\Phi\le1\) supported away from zero and one,
\begin{align*}
 &\sum_{Ax+By=N}\Phi(Ax/N)D_U(x)D_V(y)\\
 &\quad=c_1c_2\Gamma_N C_{\mathcal E}J_{\mathcal E}(N)
       \frac{N}{AB L^\kappa}
       \left(\int_0^1\Phi(t)\,dt+O(L^{-1/2})\right)+O(N^{1/2}).
\end{align*}
The statement is uniform in target valuations. Endpoint strips of
relative length \(\epsilon\) have upper bound
\(O(\epsilon NJ_{\mathcal E}(N)L^{-\kappa}+N^{1/2})\).
For every interval \(I\subset[0,1]\), the unsmoothed model correlation
has the same leading coefficient multiplied by \(|I|\), with absolute
error \(O(NJ_{\mathcal E}(N)L^{-\kappa-\delta_0})\), where
\(\delta_0=1/(8192\varphi(L_0))\).
\end{proposition}
\begin{proof}
Discard records larger than \(E\) using \eqref{eq:record-tail}.
The rough logarithms on the rest are \(L+O(\sqrt L)\), giving
\[
 \frac{M_{R,K}(n_{S'})}{V_S}
 =\frac{c_H L^{\rho-1}}{\alpha_0E_-(y)}(1+O_K(L^{-1/2})).
\]
Count the remaining small-prime sieve set in the fixed residues.
The fundamental lemma applies at level \(N^{\epsilon_0}\), for
small fixed \(\epsilon_0>0\), with
\(u=\log(N^{\epsilon_0})/\log y\asymp L^{3/5}\).
Its relative error is \(e^{-cu\log u}\), and CRT gives endpoint
error \(O(\nu(d))\) at each squarefree sieving modulus \(d\).
Summing those errors costs \(N^{\epsilon_0}L^{O(1)}\).
This is the standard fixed-dimensional fundamental-lemma application
\cite[Theorem 11.7]{richert}; target-root mergers
decrease its dimension and preserve uniformity.

The positive integer line is \((x,y)=(x_0+Bt,y_0-At)\), with
parameter length \(N/(AB)\). Modulo \(L_0\), this parameter bijects
onto all roots: at each prime at least one coefficient is a unit.
If \(R\) roots have the prescribed states, their raw probability is
\(R/L_0\). Multiplying the normalized marginal projections gives
\[
 \frac{R}{L_0}\frac{L_0c_1}{|H_1|}\frac{L_0c_2}{|H_2|}
       =c_1c_2\Gamma_N.
\]
A prime forbidden on only one side has relative factor one.
A prime forbidden on both sides has two excluded roots, merging
exactly when \(p\mid N\), and relative factor
\[
 \frac{1-(2-\1_{p\mid N})/p}{(1-1/p)^2}.
\]
Their product up to \(y\) is
\(C_{\mathcal E}(y)J_{\mathcal E,y}(N)\). The ordinary tail is
\(O(1/y)\), and the target-prime tail is \(O(L/(y\log y))\).
Partial summation inserts \(\Phi\), costing its total variation.
The same upper sieve on an endpoint strip gives its length times
the sparse density, with a power-smaller endpoint error.
To obtain the uniform interval assertion, put
\(\epsilon=L^{-\delta_0}\) and remove strips of this width near
zero and one. On the rest the rough logarithms are
\(L+O(\sqrt L+\delta_0\log L)\), giving the same relative
\(O(L^{-1/2})\) replacement. The fundamental-lemma count in any
remaining interval has the same CRT endpoint error; partial summation
costs the total variation, bounded independently of the interval.
Restoring the strips costs \(O(\epsilon)\) of the sparse scale.
All other tails are superlogarithmic or power-smaller, proving the claim.
\end{proof}

\subsection{Verification of the comparison hypotheses}
The indicators lie between zero and one, and the model bound after
\eqref{eq:corrected-model} gives \(D_U\le1\) eventually for each
fixed depth. Thus (C1) holds, with a threshold allowed to depend on
that depth. The full-unit-group exact model also lies between zero
and one.

Choose banks of primes \(p\equiv1\pmod{L_0}\), with small banks in
\([\exp(L^{1/8}),\exp(L^{1/4})]\) and big banks in
\([\exp(L^{1/3}),\exp(L^{2/5})]\). Delete primes dividing \(N\).
Every bank prime is a unit at \(A,B\), belongs to both allowed
sets, and preserves the fixed states. Hence
\begin{equation}\label{eq:exact-invariance}
 f_U(pn)=f_U(n),\quad D_U(pn)=D_U(n),\quad \xi_U(pn)=\xi_U(n)
\end{equation}
exactly, including higher powers. The deletion changes reciprocal
mass by at most \(O(L/(P_{\min}\log P_{\min}))\), which is smaller
than every fixed inverse logarithmic power. This verifies (C2).

For (C3), fix an \(S\)-smooth record \(r\le E\). If \(f(r)=0\),
both the indicator and model on that record vanish. Otherwise
\(r\bmod L_0\in H\), and insert the residual state \(U/r\)
into the periodic test of Lemma~\ref{lem:rough-means}. Its smooth
scale is \(Y/r\), so the error is \((Y/r)L^{-D}\), as required.
An additional periodic test of period at most \(L^B\) increases the
period by only a fixed factor. For every required accuracy, the
Hankel remainder works for every sufficiently large fixed depth;
constants and thresholds, rather than earlier graph exponents, may
depend on that depth. This proves (C3) with all surviving character
branches retained.

It remains to check the positive envelope (C4). Write
\(\mu_{\mathcal B}\) for the law \eqref{eq:bank-law}.

\begin{lemma}\label{lem:envelope}
For a nonnegative record function \(F(r,s)\le L^B\), and either
indicator or model on each side, its weighted correlation on the
positive line restricted by \(Ax\in J\subset[0,N]\) is at most
\[
 C\frac{|J|J_{\mathcal E}(N)}{AB L^\kappa}\E_{\mu_{\mathcal B}}F
       +NL^{O(1)}e^{-cL^{1/10}}+O(N^{1/2}).
\]
The constant may depend on the fixed coefficients.
\end{lemma}
\begin{proof}
Discard records \(r,s>E\) using \eqref{eq:record-tail} and bounded
weights. There are at most \(E^2\) retained record pairs. A shared
bank prime makes the stratum empty, because no bank prime divides
\(N\). For coprime records, \(r\mid x,s\mid y\) select one
progression of length \(|J|/(ABrs)\). Exactness excludes two roots
at an unselected bank prime and the remaining variable's root at
a prime dividing \(rs\). Thus the bank probability is exactly
\[
 \frac1{rs}\prod_{\substack{p\in\mathcal B\\p\nmid rs}}(1-2/p)
             \prod_{p\mid rs}(1-1/p)=\mu_{\mathcal B}(r,s).
\]
Forbidden-prime slopes remain units. Apply the upper Selberg sieve
with \(z=N^{1/10}\), using primes up to \(w=z^{1/16}\), which
contains all banks for large \(N\). The squared-weight CRT remainder
is \(O(z^{4+\eta})\), for arbitrarily small fixed \(\eta>0\),
uniformly in the records. After summation it is
\(N^{0.45+o(1)}=O(N^{1/2})\).

For two indicators the nonbank sieve denominator has size at least
a constant times \((\log w)^\kappa/J_{\mathcal E}(N)\).
For model times indicator it has size
\((\log y)^{1-\rho_1}(\log w)^{1-\rho_2}/J_{\mathcal E}(N)\);
the pointwise model bound \(L^{-3(1-\rho_1)/5}\) restores precisely
\(L^{-\kappa}\). For two models the denominator is
\((\log y)^\kappa/J_{\mathcal E}(N)\), and the two model bounds
give the same result. At a prime forbidden on one side only that
side's denominator occurs. Target mergers cost at most the full
factor \(J_{\mathcal E}(N)\). Sum the retained main bounds against
\(F\) to obtain the expectation.
\end{proof}

\begin{proof}[Proof of Theorem~\ref{thm:general}]
The preceding checks verify (C1)--(C4), with
\(\mathcal S_N=NJ_{\mathcal E}(N)/(AB L^\kappa)\).
This scale is between \(NL^{-C}\) and \(NL^C\) for a fixed \(C\),
by \(J_{\mathcal E}\ge1\) and Mertens' upper product bound.
Theorem~\ref{thm:model-comparison} replaces both indicators by their
models with error \(O(\mathcal S_NL^{-\delta_0})\), uniformly in
the interval. Proposition~\ref{prop:model-correlation} evaluates that
model count with exactly the coefficient in the theorem.
If \(\Gamma_N=0\), the residue restrictions give zero identically.
Otherwise its finitely many positive values have a fixed positive
minimum, yielding the relative assertion.
\end{proof}


\section{The shared graph comparison and reflected target grids}\label{sec:graph}

This section proves the common graph argument for
Theorem~\ref{thm:model-comparison}, then the residue-energy estimate
needed for its reflected family. The fixed-shift major arcs are treated
separately in Section~\ref{sec:fixed-shift-comparison}. No arithmetic
form of the indicators is used here beyond \textup{(C1)--(C3)}.
Two structural points matter. The labels lie in a fixed reduced class,
so the character frame must identify restrictions differing by a known
unit phase. A target-sized displacement requires retaining its exact
rational grids in the energy estimate.

Write \(L=\log x\).  There are \(\asymp\log L\) small banks and big
banks, with primes respectively in
\[
 [\exp(L^{1/8}),\exp(L^{1/4})],\qquad
 [\exp(L^{1/3}),\exp(L^{2/5})].
\]
Every label is \(c\bmod L_0\); primes dividing the nonzero displacement
\(k\) are removed. We use \(k=T=x\) for sums and fixed \(k=h\)
for shifts. All banks are disjoint. Their laws are
\(\mu_g(p)=(V_gp)^{-1}\), where \(V_g=\sum_{p\in\mathcal P_g}p^{-1}\)
is bounded above and below by positive fixed constants.
Banks in this section may be any active subfamilies of the banks
constructed from the class in \textup{(C2)}, provided that each size has
at least a fixed positive multiple of \(\log L\) banks.
All constants may depend on these fixed bounds and on \(L_0\).

\subsection{The character frame for progression banks}

The ideal-kernel construction in \cite[Section~4]{weightedgraphs}
uses primitive characters of conductor at most \(Q=L^{q_*}\).
On labels \(c\bmod L_0\), this list is redundant up to phases.
For example, the principal character and \(\chi_{-4}\) agree on
primes \(1\bmod4\) and are opposite on primes \(3\bmod4\).
The quotient below removes this redundancy while retaining its phase.

Let \(\mathcal X_Q\) be the primitive characters of conductor at most
\(Q\), and define
\[
 \eta\sim\eta'
 \quad\Longleftrightarrow\quad
 \operatorname{cond}(\eta\overline{\eta'})\mid L_0.
 \tag{5.1}\label{eq:gr-equivalence}
\]
The conductor here belongs to the primitive character inducing the
quotient on a common modulus.  Choose one representative from each
class and denote the resulting set by \(\mathcal R_Q\).
Thus \(|\mathcal R_Q|\le Q^2\).

To verify the equivalence and its meaning, put
\(M_0=\lcm(L_0,\operatorname{cond}\eta,\operatorname{cond}\eta')\).
Reduction from \((\Z/M_0\Z)^\times\) onto
\((\Z/L_0\Z)^\times\) is surjective.  Its kernel consists of the
units \(1\bmod L_0\).  A character is trivial on that kernel exactly
when it factors through the latter group.  This is equivalent to
its primitive conductor dividing \(L_0\). Consequently, if
\(\eta\sim\eta'\) and their quotient is induced by
\(\vartheta\bmod L_0\), then on a product of \(r\) labels,
\[
                 \eta(D)=\vartheta(c)^r\eta'(D).
\]
Both comparison products have the same slot type and the same phase.
Inequivalent characters have zero average on the kernel, and also
on any fixed coset of it: translation multiplies their sum by a unit phase.

Fix a nonempty set \(\Lambda\) of at most \(m\) slots per big bank,
where \(m\) is fixed before the eventual graph mark length.
Let \(D_\Lambda\) be the product of independent labels in these slots.
Partition its logarithm into cells
\(I_j=[jh,(j+1)h)\), and write
\(\nu_j=\Prob(\log D_\Lambda\in I_j)\).
For inequivalent \(\eta,\eta'\in\mathcal R_Q\), and any prescribed
\(A>0\), Siegel--Walfisz gives
\[
 \left|\E\big[
  \1_{\log D_\Lambda\in I_j}
        \eta(D_\Lambda)\overline{\eta'(D_\Lambda)}\big]\right|
       \ll_A L^{-A}.                                  \tag{5.2}\label{eq:gr-cell}
\]
Here and below the accuracy in Siegel--Walfisz is chosen after all
fixed logarithmic exponents.

Indeed, freeze all but one label.  Its allowable values form an
interval inside its original bank interval, still restricted to
\(c\bmod L_0\). On a common modulus at most \(L_0Q^2\), the sum of
\(\eta\overline{\eta'}\) over the appropriate reduced residue classes
is zero.  Siegel--Walfisz, summed over these classes, therefore gives
\[
 \sum_{\substack{p\le z\\p\equiv c\ (L_0)}}
       \eta(p)\overline{\eta'(p)}
       \ll_K L_0Q^2z(\log z)^{-K}.
\]
Partial summation bounds the reciprocal-prime sum on any subinterval
of \([\exp(L^{1/3}),\exp(L^{2/5})]\) by
\(O_K(L_0Q^2(L^{1/3})^{1-K})\).
This is an absolute bound; no lower bound for the subinterval's
length or harmonic mass is needed.  Divide by \(V_g\), then average
the frozen labels.  Deleting target divisors changes the answer by at most
\[
 \sum_{\substack{p\mid k\\p\ge\exp(L^{1/3})}}\frac1p
 \le \frac{\log(2|k|)}{L^{1/3}}\exp(-L^{1/3}),               \tag{5.3}\label{eq:gr-deletion}
\]
which is smaller than every fixed inverse power of \(L\).
This proves \eqref{eq:gr-cell} uniformly for the actual banks.

On cells with \(\nu_j\ge\xi\), define
\[
 K_\Lambda(D,X)=
 \sum_{\nu_j\ge\xi}
 \frac{\1_{\log D\in I_j}\1_{\log X\in I_j}}{\nu_j}
 \sum_{\eta\in\mathcal R_Q}\overline{\eta(D)}\eta(X).
 \tag{5.4}\label{eq:gr-kernel}
\]
Both product variables have the same law.  The elementary bounds are
\[
 |K_\Lambda|\le Q^2\xi^{-1},\qquad
 \sup_X\E_D|K_\Lambda(D,X)|\le Q^2,\qquad
 \#\{j:\nu_j>0\}\ll_m1+(\log L)L^{2/5}/h.              \tag{5.5}\label{eq:gr-kernel-bounds}
\]
The last inequality uses only the range of the logarithm.

On a retained cell, the kernel is the positive frame operator
\[
 f\longmapsto
 \sum_{\eta\in\mathcal R_Q}\eta\,\langle\eta,f\rangle
\]
for the conditional product probability measure.
Its Gram matrix has diagonal entries one, since all labels exceed
\(L_0Q^2\), and off-diagonal entries
\(O_A(\xi^{-1}L^{-A})\) by \eqref{eq:gr-cell}.
Thus its norm is at most
\(1+O_A(Q^2\xi^{-1}L^{-A})\le2\).
Distinct cells are orthogonal blocks, so
\[
             \|K_\Lambda\|_{2\to2}\le2.                \tag{5.6}\label{eq:gr-frame-norm}
\]
This also holds for tuple tests: conditional expectation from tuples
to their product is a contraction.
If \(\psi\) is a character of modulus \(r\le Q\), its primitive
inducing character is equivalent to one member of \(\mathcal R_Q\).
For \(r_\Lambda\) selected slots and \(\psi=\vartheta\eta\) on
these units, that summand gives
\(\vartheta(c)^{r_\Lambda}\eta(X)=\psi(X)\). The other summands give
\[
 \E_D K_\Lambda(D,X)\psi(D)
       =\psi(X)+O_A(Q^2\xi^{-1}L^{-A})                 \tag{5.7}\label{eq:gr-reproduce}
\]
at every retained \(X\).  This is where quotienting the frame is
essential: the unquotiented frame can reproduce the same constant
several times.

\begin{lemma}[Progression ideal comparison]\label{lem:gr-ideal}
For every prescribed ideal accuracy \(E_{\rm id}\) and fixed phase
range \(|\zeta|\le L^{C_4}\), the progression banks admit the signed
comparison family of \cite[Theorem~4.1]{weightedgraphs}.
Its ideal norm is at most \(L^{-E_{\rm id}}\), uniformly in every
real additive frequency.  Its coefficients ignore small-bank labels
and shared labels, free at most \(m\) slots per big bank, and have
total supremum at most \(L^{C_1}\).
The parameters \(m,C_1\) and all kernel powers are fixed before the
eventual mark length.
\end{lemma}

\begin{proof}
We give the comparison estimates used in the construction, including
the only step where restricting the prime labels changes the proof.

For a product formed from \(k_g\le m\) draws in each of at most
\(C_0\log L\) disjoint banks, unique factorization gives its exact law
\[
 n\Prob(D_\Lambda=n)
   =\prod_g\frac{k_g!}{V_g^{k_g}\prod_{p\in\mathcal P_g}a_p!},
 \qquad
 \Prob(D_\Lambda=n)\le L^{K_0}/n,                     \tag{5.8}\label{eq:gr-productlaw}
\]
where
\(K_0=C_0(\log(m!)+m\log\max(1,v_-^{-1}))\) and
\(V_g\ge v_->0\).
The factor \(C_0\) multiplies both terms of this exponent.
For a tuple test, Cauchy--Schwarz on each product fibre gives an
integer coefficient sequence of squared norm at most
\(L^{K_0}/U\) times the tuple squared norm on a product dyad
\([U,2U)\).

The elementary integer bilinear estimate
\cite[Lemma~4.2]{weightedgraphs}, or \cite[Lemma~6.2]{walsh}, gives
\[
 \left|\sum_{u,v}a_ub_v\ee{\theta uv}\right|^2
 \ll \|a\|_2^2\|b\|_2^2
             (1+U/r)(V+r\log(2r))                  \tag{5.9}\label{eq:gr-integer-bilinear}
\]
when the two supports lie in containing intervals of lengths
\(O(U),O(V)\), \((a,r)=1\), and
\(|\theta-a/r|\le r^{-2}\).
For completeness, Cauchy--Schwarz and expansion leave
\(V+\sum_{1\le j\ll U}\min(V,\|\theta j\|^{-1})\).
Divide the \(j\)'s into blocks of diameter less than \(r/2\).
Within a block the rational parts are \(1/r\) separated and their
perturbations differ by at most \(1/(2r)\); hence the sum is
\(O(V+r\log(2r))\) per block.  This proves
\eqref{eq:gr-integer-bilinear} and permits arbitrary subsets of the
containing intervals.

Apply this estimate to two nonempty selected product sets \(A,B\).
On product dyads \(U,V\), use Dirichlet approximation with denominator
ceiling \(\lfloor UV/S\rfloor\).  When the resulting denominator
exceeds \(Q\), \eqref{eq:gr-productlaw} gives the normalized bound
\[
 \ll L^{K_0}\big(o(L^{-A})+Q^{-1}+L^{7/5}/S\big)^{1/2}.
 \tag{5.10}\label{eq:gr-highden}
\]
There are at most \(L^{7/5}\) product dyads per variable.
The same bound applies to the kernel terms by moving each kernel
onto its endpoint test using \eqref{eq:gr-frame-norm}.
The factor \((uv)^{i\zeta}\) separates and preserves test norms.

For denominators \(r\le Q\), all products are units modulo \(r\).
The multiplicative-character expansion of the rational phase has
total absolute coefficient at most \(\sqrt{\varphi(r)}\).
Replace each primitive character by its representative in
\(\mathcal R_Q\) and apply \eqref{eq:gr-reproduce}.
Grouping characters after inserting their unit phases cannot increase
coefficient mass, so the reproduction error is
\(O_A(Q^{5/2}\xi^{-1}L^{-A})\).
Write the remaining phase as
\(M(v)=\exp(i\zeta v)\ee{\beta\exp v}\).
On matched cells its derivative is
\(O(L^{C_4}+S)\), and the two logarithmic products move by at most
\(2h\).  By \eqref{eq:gr-kernel-bounds}, this costs
\(O(Q^4h(S+L^{C_4}))\).
The union of discarded cells has probability at most
\((N_A+N_B)\xi\), where \(N_A,N_B\) count the possible cells.
Cauchy--Schwarz bounds its contribution by the square root of this
probability times the two test norms.
These estimates prove the two-product comparison, for arbitrary
\(L^2\) tuple tests, to any prescribed power \(L^{-A_0}\).

There is no circular choice of precisions.  One possible hierarchy is
\[
 \begin{split}
 B_0&>A_0+2(7/5)+10,\\
 q_*,s_*&>2(B_0+K_0+10)+7/5,\\
 h_*&>4q_*+\max(s_*,C_4)+B_0+10,\\
 x_*&>2B_0+7/5+h_*+10,
 \end{split}
\]
with \(Q=L^{q_*},S=L^{s_*},h=L^{-h_*},\xi=L^{-x_*}\).
Choose the accuracy in \eqref{eq:gr-cell} afterwards.
All these choices are independent of the eventual mark length.

Here is the signed cancellation which completes the ideal construction.
Split the big banks into two blocks of comparable size and designate
\(m\) probe slots per group.  For nonempty probe subsets
\(A,B\) in the two blocks, free their slots and give the comparison
coefficient
\[
              -(-1)^{|A|+|B|}K_A(D_A,X_A)K_B(D_B,X_B).
 \tag{5.11}\label{eq:gr-signed}
\]
Here \(D_A,D_B\) are auxiliary-label products and \(X_A,X_B\)
are the unshared endpoint products.  The coefficient has no small
or shared labels.  There are at most \(2^{mC_0\log L}\) comparisons;
by \eqref{eq:gr-kernel-bounds} their total supremum is at most
\[
        1+2^{mC_0\log L}Q^4\xi^{-2}\le L^{C_1}.         \tag{5.12}\label{eq:gr-coeffcost}
\]
Expanding both kernels also separates each coefficient into a sum
of products of four bounded factors: one on each auxiliary product
and one on each unshared endpoint list.  The number of terms and
total scalar coefficient mass are fixed powers of \(L\).

We now make the exceptional projection and the cancellation explicit.
Write the two big-bank blocks as \(\mathcal B_1,\mathcal B_2\), of
sizes \(s_1,s_2\), and their designated probe coordinates as
\(\mathcal I_1,\mathcal I_2\).  For a coordinate \(i\), let \(E_i\)
average its independent label.  These commuting operators are
orthogonal projections.  Let \(Q_{\rm tail}\) average the last
\(\ell\) input or target coordinates in every group, and let
\(\Pi\) be within-group symmetrization on all \(M=J+\ell\) coordinates.
The ideal multipliers do not use those last coordinates.
For symmetric endpoint tests \(F,H\), put
\(F'=Q_{\rm tail}F,\ H'=Q_{\rm tail}H\).
Both averaging operations are contractions.

Let \(D_{\rm sh}\) be the product of the first \(J\) coordinates,
and write
\[
 \Phi_{\Theta,\zeta}
       =D_{\rm sh}^{i\zeta}\ee{\Theta D_{\rm sh}},\qquad
 R(F',H')=\E\big[F'\overline{H'}\Phi_{\Theta,\zeta}\big].
\]
Thus \(R\) is a bilinear pairing of norm at most one.
Resolve the probes in the first block of \(F'\) and the second
block of \(H'\):
\[
 \begin{split}
 F'&=\sum_{S\subset\mathcal I_1}F_S,\qquad
 F_S=\prod_{i\in S}(I-E_i)
             \prod_{i\in\mathcal I_1\setminus S}E_iF',\\
 H'&=\sum_{T\subset\mathcal I_2}H_T,\qquad
 H_T=\prod_{i\in T}(I-E_i)
             \prod_{i\in\mathcal I_2\setminus T}E_iH'.
 \end{split}
\]
These are orthogonal decompositions.  Define
\(P_1=\prod_{i\in\mathcal I_1}(I-E_i)\),
\(P_2=\prod_{i\in\mathcal I_2}(I-E_i)\), so that
\(F_{\rm full}=P_1F'\), \(H_{\rm full}=P_2H'\).

In one group the original symmetric vector decomposes orthogonally
over subsets of its \(M\) coordinates.  Symmetry makes the squared
norms equal among subsets of cardinality \(j\).
Requiring all \(m\) probes to occur and all last \(\ell\) slots
to be absent therefore retains exactly the fraction
\[
 \frac{\binom{M-m-\ell}{j-m}}{\binom Mj}
     =\frac{(j)_m(M-j)_\ell}{(M)_{m+\ell}}
     \le 2(j/M)^m(1-j/M)^\ell
     \le2\ell^\ell m^{-\ell}=:\theta_m
\]
when \(M\ge(m+\ell)(m+\ell-1)\); impossible cardinalities give zero.
For a vector which is not a tensor product across groups, first
resolve by the tuple of cardinalities in the groups of a block.
Independent within-group permutations make all subset tuples of
those cardinalities have equal squared norms.  The retained fraction
is consequently the product of the displayed fractions.
Summing over the cardinality tuples proves the operator estimates
\begin{equation}\label{eq:gr-blockprojection}
 \|P_1Q_{\rm tail}\Pi\|\le\theta_m^{s_1/2},\qquad
 \|P_2Q_{\rm tail}\Pi\|\le\theta_m^{s_2/2}.
\end{equation}
Averaging tails outside the block is a further contraction.
We use \(\ell=1\), and each \(s_i\ge c_0\log L/3\).

For a component pair \(F_S,H_T\), put
\(U=\mathcal I_1\setminus S\), \(V=\mathcal I_2\setminus T\).
A comparison freeing \(A\subset\mathcal I_1\),
\(B\subset\mathcal I_2\) vanishes unless \(A\subset U,B\subset V\).
Indeed, if a freed input coordinate belongs to \(S\), it is
mean zero, absent from the target, and absent from the comparison
multiplier: integrating it gives zero.  The corresponding statement
holds for a freed target coordinate belonging to \(T\).
This uses the orientation in \eqref{eq:gr-signed}: the \(A\)
endpoint labels are from the target and the \(B\) endpoint labels
are from the input.

If \(A\subset U,B\subset V\), condition on all shared coordinates
outside \(A\cup B\), with product \(C\).
The input test is independent of the freed \(A\) labels and the
target test is independent of the freed \(B\) labels.
The two-product kernel comparison established above therefore
replaces the unsigned comparison by \(R(F_S,H_T)\), with error
\[
                    L^{-A_0}\|F_S\|_2\|H_T\|_2.
\]
Its frequency is \(\Theta C\); uniformity for every real frequency
is essential.  Cauchy--Schwarz over the conditioned coordinates
gives exactly the displayed unconditional norms.

For nonempty \(U,V\), the total unsigned subset sign is
\[
 \sum_{\substack{\varnothing\ne A\subset U\\
                 \varnothing\ne B\subset V}}
       (-1)^{|A|+|B|}
 =\left(\sum_{\varnothing\ne A\subset U}(-1)^{|A|}\right)
  \left(\sum_{\varnothing\ne B\subset V}(-1)^{|B|}\right)=1.
\]
The initial minus sign in \eqref{eq:gr-signed} cancels the raw
pairing.  If either missing set is empty, no comparison survives.
All exceptional raw component pairs together have the exact identity
\[
 R_{\rm exc}(F,H)=
 R(P_1F',H')+R((I-P_1)F',P_2H').
\]
In operator notation, if \(M_\Phi\) multiplies by
\(\Phi_{\Theta,\zeta}\), their operator on the symmetric space is
\[
 T_{\rm exc}=
 \Pi Q_{\rm tail}
   \big[M_\Phi P_1+P_2M_\Phi(I-P_1)\big]
                 Q_{\rm tail}\Pi .
\]
Consequently
\[
 \|T_{\rm exc}\|
       \le\theta_m^{s_1/2}+\theta_m^{s_2/2}.
\]
This identity sums the exceptional components before taking an
absolute value.  It is the reason their bound has no subset-count
or comparison-coefficient factor.

There are at most \(3^{P}\) admissible component/free-subset triples,
where \(P=m(s_1+s_2)\le mC_0\log L\): each probe is indicated,
missing but not freed, or freed.  Thus the full signed ideal
operator satisfies
\[
 \|T_{\rm signed}-T_{\rm exc}\|
           \le3^{P}L^{-A_0}.
\]
Choose \(m\) so that \(\theta_m<1\) and
\(-c_0\log\theta_m/6>E_{\rm id}+3\), and then choose
\(A_0>E_{\rm id}+mC_0\log3+10\).
The kernel precisions above and the exponent \(C_1\) in
\eqref{eq:gr-coeffcost} are chosen after these choices.
Only afterwards require \(J\ge m\) and
\(J+\ell\ge(m+\ell)(m+\ell-1)\); increasing this fixed \(J\)
to meet the physical root-moment theorem changes no kernel
parameter or coefficient exponent.
The resulting ideal norm is at most \(L^{-E_{\rm id}}\).
This proves the asserted ideal norm and coefficient conditions.
\end{proof}

\subsection{A large displacement in the physical graph}

\begin{lemma}[Large-target root moment and pairing]\label{lem:gr-large}
The root moment and endpoint pairing from
\cite[Theorem~3.5 and Corollary~3.11]{weightedgraphs} hold for any
nonzero integer \(k\) coprime to the bank primes, with the same
logarithmic accuracies.  In particular they hold for \(k=T\asymp x\)
after target-prime deletion.  The first endpoint is required to be
independent of its marks, exactly as in the imported pairing.
\end{lemma}

\begin{proof}
Fix the labels, patterns, and permutations of a closed walk of length
\(2R\), where \(R=\lceil L^{0.52}\rceil\).
Its position offsets are \(ks_i\).  Return to the initial integer
is equivalent to \(s_{2R}=0\), independent of the nonzero \(k\).
The root dependence consists of divisibility and damping conditions
at the offsets \(ks_i\).

The independent-residue replacement
\cite[Lemma~3.6]{weightedgraphs} is uniform in the numerical values
of these offsets.  Indeed its active modulus has logarithm
\(O_M(R(\log L)L^{2/5})=o(L)\).
The Bonferroni truncation has the same modulus and term count,
and each interval CRT error is \(O(1)\), independently of the
specified residues.  Once the deterministic damping from active
primes is factored out, the error is relative and can be summed with
absolute path coefficients using the imported absolute row bound.

In the resulting independent residue space, replace the root residue
\(z_p\) by \(k^{-1}z_p\bmod p\) for every bank prime \(p\).
Every fixed-path divisibility and damping expectation is then exactly
the expectation for \(k=1\).
The remaining signed coefficients are unchanged.
Thus the \(k=1\) root moment gives the identical bound for this \(k\).
Negative root intervals follow either from the same CRT count or
from \(n\mapsto-n,\ k\mapsto-k\).

It remains to avoid a large factor in the endpoint pairing.
Let \(D_{\max}\) bound the edge products, so that
\(\log D_{\max}=O_M((\log L)L^{2/5})\).
Partition the first-endpoint interval \(I\) into blocks of length
\(\lceil |k|D_{\max}\rceil\), as in the source pairing, and refine
each block by residue modulo \(|k|\).
The operator preserves these residues.
Each refined block has \(O(D_{\max}+1)\) integer positions.
For \(B=\mathcal A_k\mathcal A_k^*\), the matrix
\((\langle u_n,B^Ru_m\rangle)\) on one such block is positive
semidefinite.  Its largest eigenvalue is at most its trace.
Here \(u_n\) is the unnormalized vector equal to one on all marks at
\(n\).  Mark independence of the first endpoint \(f_n\) gives
\[
 \langle f_j,B^Rf_j\rangle
 \le \Big(\sum_{n\in I_j}|f_n|^2\Big)
             \sum_{n\in I_j}\langle u_n,B^Ru_n\rangle.
\]
Spectral H\"older and Cauchy--Schwarz give the source block estimate;
its only dimension factor is now
\[
 \big(O(D_{\max})\sup|f_n|^2\big)^{1/(2R)}
 \le
 \exp\!\left(O_M\!\left(
 \frac{(\log L)L^{2/5}+\sqrt L}{L^{0.52}}\right)\right)=O(1).
 \tag{5.13}\label{eq:gr-block}
\]
Second-endpoint neighborhoods have bounded overlap within a residue
class and are disjoint between residue classes.
H\"older over blocks uses only the sum of root traces on the whole
original interval, bounded by the root theorem.
No CRT estimate on an individual residue class modulo \(k\) is used.
This proves the pairing uniformly in \(k\).
\end{proof}

For bounded invariant endpoint functions, the list moments needed
here do not create a mark-dependent logarithmic loss.  For a fixed
ordered distinct list with product \(P\), divisibility gives
\(O(X/P)\) integers in an interval of length \(O(X)\) contained in
\([-CX,CX]\).  For fixed mark length, \(P=x^{o(1)}\).
After summing the list weights,
\[
 \sum_{\text{lists}}\frac1P\prod_gV_g^{-M}\le1.
 \tag{5.14}\label{eq:gr-list-moments}
\]
Thus the endpoint \(r\)-moments are \(O_r(X)\).
Condition \textup{(C1)} bounds every fixed-depth model without a
depth-dependent logarithmic exponent. Thus this observation is valid
before the final depth is chosen.

\subsection{Shared labels and the two comparison geometries}

For sums take \(k=x\); for shifts take fixed \(k=h\ne0\).
We first obtain the same comparison identity in both cases.

We spell out the arithmetic consequence of the pairing, since
the large reflected target is not covered by the fixed-shift theorem
of \cite{walsh}.
For a vector \(\mathbf t=(t_g)\) and a bounded function \(b\) on its
ordered distinct lists, put
\[
 W_{\mathbf t}^b(n)=
 \prod_g\big(\1_{\omega_g(n)=t_g}V_g^{-t_g}\big)
       \sum_{\substack{\mathbf p\text{ a }\mathbf t\text{-list}\\
                        p_{g,j}\mid n}} b(\mathbf p).
 \tag{5.15}\label{eq:gr-listweight}
\]
Here \(\omega_g(n)\) counts distinct bank-prime divisors of \(|n|\).
For fixed \(t_g\le m+1\), these weights are bounded by a fixed
power of \(L\).  In particular \(W_{\mathbf1}^1=\Omega_{\mathcal P}\).

Let \(\alpha,\beta\) be bounded functions invariant under multiplying
by every bank prime. For sums extend the first centered function to
the negative half-line by absolute value and the second function to
the positive half-line. For shifts use positive endpoints on both
sides. These sign restrictions preserve prime-multiplication invariance.
Take a smooth \(\Psi\) supported in a compact subset of \((0,\infty)\).
The pairing and Lemma~\ref{lem:gr-ideal} imply
\[
 C_0+\sum_\nu c_\nu C_\nu=O_A(L^{-A}),\qquad
 \#\{\nu\}+\sum_\nu|c_\nu|\le L^{C},                  \tag{5.16}\label{eq:gr-comparison}
\]
where
\[
 C_0=\sum_{w\ne0}\frac{\Psi(|w|/x)}{|w|}
       \alpha(w)\beta(w+k)
       W_{\mathbf1}^1(w)W_{\mathbf1}^1(w+k)
\]
and each comparison has the form
\[
 \begin{split}
 C_\nu=\E_{D_1,D_2}a_\nu(D_1)b_\nu(D_2)
 \sum_{w\ne0}\frac{\Psi(|w|/(xD_1D_2))}{|w|}
  \alpha(w)\beta(w+kD_1D_2)\\
 {}\times W_{\mathbf t}^{u_\nu}(w)
                  W_{\mathbf t}^{v_\nu}(w+kD_1D_2).
 \end{split}                                                     \tag{5.17}\label{eq:gr-comparison-term}
\]
The two auxiliary products are independent, use at most \(m\) draws
per big bank, and each contains at least one big-bank prime.
The four separated coefficient functions have modulus at most one.
The list coefficients \(u_\nu,v_\nu\) ignore all small-bank labels;
in every small bank \(t_g=1\).

Here is the derivation of these precise assertions from the physical
graph.  A pattern has \(M-1\) dilation slots per bank, some shared
and some freed, and one distinguished unshared endpoint slot.
Write its edge product as \(D=D_RD_C\), with shared product \(D_R\)
and auxiliary product \(D_C\).
For a source of either sign in \(|n|\in[Z,2Z)\), an originally nonzero edge
weight satisfies
\(\Psi(|n|/(xD))\ne0\), hence \(D\ll Z/x\).
Since \(|k|\ll x\), this implies \(|n+kD|\ll Z\).
Insert a second-endpoint cutoff on \([-C_\Psi Z,C_\Psi Z]\)
before Mellin inversion.  It equals one on every contributing edge.
After inversion retain this same cutoff.  An edge originally having
zero \(\Psi\) still has zero full Mellin integral, so the insertion
does not change the identity.  No cutoff involving shared-label
products has been introduced.

Mellin inversion of \(\Psi\), with the harmonic weight \(1/|n|\),
now applies the pairing of Lemma~\ref{lem:gr-large} on each of
\(O(L)\) source dyads.  Its tails are bounded by the absolute
row/column estimate, its main part by the small ideal norm, and
its endpoint norms by \eqref{eq:gr-list-moments}.
Compressing the physical operator to exact-count vertices
\(\omega_g(n)=\omega_g(n')=M\) is legitimate:
the corresponding projection depends only on the integer, preserves
mark independence, and sets both damping factors equal to one.

For completeness, put \(n=D_Rw\), \(n'=D_Rw'\), so
\(w'=w+kD_C\), and put \(t_g=1+\#\{\text{freed slots in }g\}\).
Invariance and harmonic scaling give
\[
 \alpha(D_Rw)=\alpha(w),\quad \beta(D_Rw')=\beta(w'),\quad
 \frac{\Psi(|n|/(xD))}{|n|}
       =\frac1{D_R}\frac{\Psi(|w|/(xD_C))}{|w|}.
\]
The graph normalization at a shared prime combines with \(1/D_R\)
to give its probability law \(\mu_g\).
More explicitly, with \((D_R)_g\) its product in bank \(g\),
\[
 \frac{V_g^{-M-t_g}}{(D_R)_g}
       =V_g^{-2t_g}
           \prod_{\text{shared slots in }g}\mu_g(p).   \tag{5.18}\label{eq:gr-sharednormal}
\]
For fixed reduced endpoints and fixed auxiliary labels,
the shared-label probability is
\[
 \prod_g\1_{\omega_g(w)=\omega_g(w')=t_g}
              +O(L^{C_M}e^{-L^{1/8}}).               \tag{5.19}\label{eq:gr-sharedprob}
\]
Indeed there are only \(O_M(\log L)\) draws, each atom is
\(O(e^{-L^{1/8}})\), and \(ww'D_C\) has \(O(L)\) prime divisors.
If the reduced counts are \(t_g\), distinct draws avoiding those
divisors reconstruct the exact graph edge.  If a reduced count is
larger, a valid reconstruction must choose a shared prime already
dividing that endpoint, an event of the displayed small probability.
This proves \eqref{eq:gr-sharedprob}, also for negative \(w\).

Summing its absolute error costs only fixed logarithmic powers,
by bounded endpoints and the harmonic weight.
The restrictions excluding auxiliary labels from endpoint lists
can then be removed: for an auxiliary prime \(p\mid D_C\),
\(p\mid w\) if and only if \(p\mid w+kD_C\); these positions have
density \(O(1/p)\) in the containing interval.
Every auxiliary prime is at least \(\exp(L^{1/3})\).
An unshared prime common to the two endpoints divides \(kD_C\);
bank-prime deletion implies it divides \(D_C\), so the same estimate
covers it.  Thus these errors remain smaller than every logarithmic
power.  Finally expand the separated coefficients from
\eqref{eq:gr-signed}.  Formula \eqref{eq:gr-sharednormal} yields
exactly \eqref{eq:gr-listweight}, while \(D_C=D_1D_2\).
The raw pattern has \(D_C=1,\mathbf t=\mathbf1\).
This proves \eqref{eq:gr-comparison}--\eqref{eq:gr-comparison-term}
with their stated independence properties.

\subsection{Prime sums at target-sized conductor}

Only this subsection uses the zero-free input.
It is important that its conductor range is polynomial in \(x\),
rather than a power of \(\log x\).

\begin{lemma}[Prime factor estimate]\label{lem:gr-primes}
Assume the fixed all-conductor zero-free half-plane stated in
Section~\ref{sec:inputs}.
Fix \(\sigma_1\in(11/12,1)\).
For \(q\le x^B\), every character \(\chi\bmod q\), and
\(P\ge\exp(L^a)\), uniformly on subintervals \(I\subset[P,2P]\),
\[
 \sum_{p\in I}\chi(p)p^{it}
 =\delta_\chi\int_I\frac{u^{it}}{\log u}\,du
       +O((1+|t|)P^{\sigma_1}L^C).                  \tag{5.20}\label{eq:gr-primes}
\]
Here \(\delta_\chi=1\) for a principal character and zero otherwise,
and the estimate holds for every real \(t\).
For progression banks the main vanishes unless \(\chi\) is induced
from a character modulo \(L_0\).
Removing target or modulus divisors preserves the error bound.
\end{lemma}

\begin{proof}
A fixed distance to the right of the zero-free boundary, the local
logarithmic-derivative formula and functional equation give
\[
 \left|\frac{L'}{L}(s,\chi)+\frac{\delta_\chi}{s-1}\right|
                 \ll \log(q(|\Im s|+3)).
 \tag{5.21}\label{eq:gr-logderivative}
\]
Imprimitive Euler factors contribute \(O(\log q)\), since the real
part is fixed and positive.
Apply truncated Perron to
\(-L'/L(s,\chi)\), at
\(c=1+1/\log Z\), height \(Z^2\), and a half-integer endpoint
\(Z\in[P/2,4P]\).
The truncation error is \(O(\log^2 Z)\): near \(Z\), sum the bounds
\(O(Z\log Z/(Z^2|n-Z|))\); away from \(Z\), use absolute convergence
at \(c\).  Shift to \(\sigma_1\).
The principal pole gives \(\delta_\chi Z\), the new vertical integral
is \(O(Z^{\sigma_1}\log^2(qZ))\), and horizontal edges cost
\(O(Z^c\log(qZ)/(Z^2\log Z))\).
Prime powers cost \(O(Z^{1/2}\log^2 Z)\).
Replacing a real endpoint by the adjacent half-integer changes its
prime-weighted count by \(O(\log Z)\).  Therefore
\[
       \sum_{p\le Z}\chi(p)\log p
          =\delta_\chi Z+O(Z^{\sigma_1}L^C).
\]
Partial summation against \(u^{it}/\log u\) proves
\eqref{eq:gr-primes} with its linear \(1+|t|\) growth.

Expand the progression condition into the finitely many characters
modulo \(L_0\).  A constituent is principal exactly when the original
character is induced from \(L_0\).
There are at most \(O(\log(q(2+|k|)))=O(L)\) deleted prime divisors;
their absolute unweighted contribution is absorbed by the error.
\end{proof}

In particular, for a nonexceptional character the prime factor saves
\(P^{\sigma_1-1}\) against its trivial length.  The all-height form
above is convenient: smooth product-cutoff Mellin transforms have
integrable \((1+|t|)\)-weighted absolute value bounded by a fixed
power of \(L\), so there is no boundary-smoothing loss from a sharp bank.

\subsection{The exact target-grid energy}

Let \(S=\{p:p\le y\}\), \(y=\exp(L^{2/5})\), and
\(E_0=\exp(\sqrt L)\), as in Section~\ref{sec:comparison}.
Let \(\xi\) be either centered function in \textup{(C3)}, at a
fixed depth which will be chosen after the required precision.
For a comparison term, define
\[
 F(n)=\psi(n/Y)n^{i\sigma}\xi(n)W_{\mathbf t}^b(n),
 \quad Y=x^{1+o(1)},\quad |\sigma|\le L^C,             \tag{5.22}\label{eq:gr-energy-F}
\]
where the support of \(\psi\) is a fixed compact positive interval.
All its logarithmic derivatives are bounded by fixed powers of \(L\).
The coefficient \(b\) ignores a selected small-bank label, whose count
is exactly one.  This holds for every comparison above.
Reflection changes only fixed character factors and additive signs,
so we state the estimate at a positive endpoint.

\begin{proposition}[Target-grid residue energy]\label{prop:gr-energy}
Suppose that
\[
       q\le xL^B,\qquad
       \frac{Y}{qP_{\max}}\gg_A L^A
                       \quad\text{for every fixed }A,
 \tag{5.23}\label{eq:gr-energy-ratio}
\]
where \(P_{\max}\le\exp(CL^{1/4})\) is the selected small-bank scale.
For every prescribed \(D>0\), a fixed sufficiently deep density model
gives
\[
 E_q(F):=\sum_{a\bmod q}\left|
           \sum_{n\equiv a\ (q)}F(n)\right|^2
                      \ll_D\frac{Y^2}{q}L^{-D}.       \tag{5.24}\label{eq:gr-energy}
\]
The estimate is uniform in all valuations of \(q\).
It also holds with \(F(n)\) multiplied by \(\ee{\beta n}\)
when \(|\beta|Y\) is bounded by a fixed power of \(L\).
\end{proposition}

\begin{proof}
Refine to \(Q=\lcm(q,L_0)\).
Cauchy--Schwarz on the refined classes gives
\(E_q(F)\le(Q/q)E_Q(F)\le L_0E_Q(F)\).
For every gcd profile \(d\mid Q\), put \(r=Q/d\).
Multiplicative orthogonality gives the exact identity
\[
 E_Q(F)=\sum_{d\mid Q}\frac1{\varphi(r)}
       \sum_{\chi\bmod r}\left|\sum_m F(dm)\chi(m)\right|^2.
 \tag{5.25}\label{eq:gr-gcdidentity}
\]
The character zeros impose \((m,r)=1\), so this identity retains
prime powers, including primes dividing \(L_0\). We do not assume
that the endpoint is prime to \(L_0\). Call a character exceptional
when its primitive inducing character has conductor dividing \(L_0\).
There are only a fixed number of such characters for each \(r\).
In particular \(L_0\mid r\) is not required.

We first identify the negligible sets in the norm actually needed.
For a set \(\mathcal B\) and \(|G|\le L^C\), Cauchy--Schwarz in each
residue class gives
\[
 E_Q(G\1_{\mathcal B})
 \le(1+CY/Q)\sum_{n\asymp Y}|G(n)|^2\1_{\mathcal B}(n).
 \tag{5.26}\label{eq:gr-tail-energy}
\]
The factor is \(O(Y/Q)\) here.
The following populations are exponentially small relative to \(Y\),
up to fixed logarithmic factors:
\[
 \begin{array}{ll}
 n_S>E_0:& O(YL^C e^{-cL^{1/10}}),\\
 p\mid n,\ p\mid Q,\ p>y:& O(YL/y),\\
 p\mid n,\ p\mid Q,\ p\text{ in the selected bank}:&
                  O(YL/P_{\min}),\\
 p^2\mid n,\ p\text{ in the selected bank}:&
                  O(Y/P_{\min}).
 \end{array}                                                       \tag{5.27}\label{eq:gr-tailsets}
\]
For the first bound, take \(a=c/\log y\). Counting multiples of
each possible smooth record gives
\[
 \#\{n\le CY:n_S>E_0\}
 \le CY E_0^{-a}\prod_{p\le y}(1-p^{a-1})^{-1}
                 \ll YL^C e^{-cL^{1/10}}.
\]
This estimate uses no property of the indicators. The other bounds
follow by counting multiples,
using \(P_{\min}\ge\exp(cL^{1/8})\).
Thus all these sets are negligible in \eqref{eq:gr-tail-energy}.

For nonexceptional characters remove only selected primes dividing
\(Q\) and selected-prime squares.  Do not yet impose \(n_S\le E_0\),
as that would couple the extracted prime to its cofactor.
The unique selected small-bank prime gives \(n=dpm\).
On a dyadic prime interval \(p\asymp P\), exact dilation invariance
and the independence of \(b\) from this slot make its remaining
coefficient \(\gamma_d(m)\) independent of \(p\), with
\(|\gamma_d(m)|\le L^C\).
It includes the condition that \(m\) has no other selected-bank prime,
and every other bank weight.  Its containing length is
\(M\asymp Y/(dP)\).
Mellin inversion of the product cutoff separates the character sum as
an integral of
\[
 A_\chi(t)B_{d,\chi}(t),\qquad
 A_\chi(t)=\sum_{p\asymp P,\ p\in\mathcal P_g}\chi(p)p^{it},
 \quad
 B_{d,\chi}(t)=\sum_{m\asymp M}\gamma_d(m)\chi(m)m^{it}.
 \tag{5.28}\label{eq:gr-primefactorization}
\]
The omitted Mellin scaling has modulus one.
The transform has polynomial-logarithmic weighted \(L^1\) norm,
including weight \(1+|t|\).
Condition \eqref{eq:gr-energy-ratio} gives
\[
                       M/r\asymp Y/(QP)\gg1,
\]
independently of \(d\).
Elementary character orthogonality therefore yields
\[
 \frac1{\varphi(r)}\sum_{\chi\bmod r}|B_{d,\chi}(t)|^2
 \le(1+CM/r)\sum_{m\asymp M}|\gamma_d(m)|^2
                  \ll \frac{M^2}{r}L^C.             \tag{5.29}\label{eq:gr-cofactorenergy}
\]
Apply Lemma~\ref{lem:gr-primes} to \(A_\chi\) and Minkowski to the
integral in \eqref{eq:gr-primefactorization}.
Its nonexceptional energy at profile \(d\) is
\[
 \ll \frac{Y^2}{Qd}L^C\exp(-cL^{1/8}).               \tag{5.30}\label{eq:gr-nonexceptional}
\]
The prime dyadic decomposition costs only a fixed power of \(L\).
Summing \(d\) is harmless because
\(\sum_{d\mid Q}1/d\le Q/\varphi(Q)\ll\log\log(3Q)\).
There is no number-of-characters or unweighted divisor-count loss.

For exceptional characters use the further restriction
\(n_S=A\le E_0\), and remove \(Q\)-prime divisors above \(y\).
The errors were already controlled in energy.
Now the gcd profile is determined by the smooth record:
\(\gcd(A,Q)=d\), \(n=Av\), \((v,S)=1\).
If \(\chi^*\) is the primitive character inducing an exceptional
\(\chi\bmod r\), then on the retained profile
\[
             \chi(n/d)=\chi^*(A/d)\chi^*(v).
\]
Indeed \(\operatorname{cond}\chi^*\mid r\) and
\((A/d,r)=1\). One must use \(\chi^*\) here: a literal character
modulo \(L_0\) could vanish on a retained bad-prime prefix.
Every other list weight depends only on \(A\), is bounded by \(L^C\),
and the remaining periodic test in \(v\) has fixed period \(L_0\).
The one-variable hypothesis must center this full periodic test;
matching only a scalar principal mean would not suffice.

There is still the condition that \(v\) avoid primes \(p\mid Q\)
above \(y\).  Remove it before applying the fixed-period centering
lemma, extending the same bounded periodic test to those new points.
At each retained record, multiples of such \(p\) number at most
\(CY/(Ap)\) on the containing interval \(0<v\le CY/A\).
Consequently the change in the profile sum is at most
\[
 L^CY\sum_{\substack{A\ S\text{-smooth}\\d\mid A}}\frac1A
                \sum_{\substack{p\mid Q\\p>y}}\frac1p
 \le \frac{Y}{d}L^CV_S^{-1}\frac{\log Q}{y\log y}
 \ll_D (Y/d)L^{-D}.                                  \tag{5.31}\label{eq:gr-avoidance}
\]
We used the exact harmonic identity
\[
 \sum_{\substack{A\ S\text{-smooth}\\d\mid A}}\frac1A
                       =\frac1{dV_S}                 \tag{5.32}\label{eq:gr-record-harmonic}
\]
when \(d\) is \(S\)-smooth; otherwise this restricted profile is empty.

After this extension, hypothesis \textup{(C3)} applies record by
record at length \(Y/A\).
It permits arbitrary bounded fixed-period tests, arbitrary choices
of retained records, and the stated smooth phases.
Choose its logarithmic accuracy after the list and cutoff costs.
Using \eqref{eq:gr-record-harmonic} gives, for every prescribed \(D_0\),
\[
            \left|\sum_mF_{\rm good}(dm)\chi(m)\right|
                         \ll (Y/d)L^{-D_0}.          \tag{5.33}\label{eq:gr-exceptional}
\]
Restoring the avoidance condition costs \eqref{eq:gr-avoidance}.
In \eqref{eq:gr-gcdidentity} the total exceptional contribution is at most
\[
 \frac{Y^2}{Q}L^{-2D_0}
           \sum_{d\mid Q}\frac1d\,\frac{Q/d}{\varphi(Q/d)}
                     \ll \frac{Y^2}{Q}L^{-2D_0+C}.   \tag{5.34}\label{eq:gr-exceptionalenergy}
\]
Choose \(D_0\) large, add the negligible sets, and use the refinement
inequality to prove \eqref{eq:gr-energy}.
Multiplying the cutoff by \(\ee{\beta Yu}\) changes only its fixed
logarithmic derivative bounds when \(|\beta|Y\le L^B\), proving the
last assertion.
\end{proof}

This proof uses only boundedness, exact prime invariance, and the
recordwise periodic means in \textup{(C1)--(C3)}. In particular no
subgroup support, pole order, positivity of an individual character
branch, or unit-support condition at the fixed modulus enters the
energy estimate.

\subsection{From target-grid energy to comparison cancellation}

We now specialize to the reflected family \(k=T=x\) and complete
its comparison estimate with all normalizations visible. The next
section gives the different frequency argument for fixed \(k=h\).
Dyadically localize the two auxiliary products in
\eqref{eq:gr-comparison-term} to \(D_i\asymp H_i\).
Then
\[
 H_i\ge\exp(L^{1/3}),\qquad
 H_i\le\exp(O_m((\log L)L^{2/5})),\qquad
 Y=xH_1H_2.
\]
Insert endpoint dyadic cutoffs before Mellin inversion, as in the
reflection argument.  The resulting ordinary comparison has form
\[
 \frac1Y\int_0^1\widehat F(\theta)\widehat G(-\theta)
                           m(T\theta)\,d\theta,\qquad
 m(\theta)=\E b_1(D_1)b_2(D_2)\ee{\theta D_1D_2}.
 \tag{5.35}\label{eq:gr-fouriercomparison}
\]
External Mellin transforms and their fixed logarithmic costs are
understood here.  The endpoint \(F\) is centered and has exactly the
form in \eqref{eq:gr-energy-F}; \(G\) has squared norm at most
\(YL^C\).  This follows from bounded invariant endpoints and the
pointwise list bound, without any multiplicativity of list coefficients.

By \eqref{eq:gr-productlaw} and
\eqref{eq:gr-integer-bilinear}, for a rational approximation
\(|\theta-a/r|\le r^{-2}\),
\[
 |m(\theta)|^2\ll L^{C_{\rm prod}}
 \left(r^{-1}+H_1^{-1}+\frac{\log(2r)}{H_2}
                    +\frac{r\log(2r)}{H_1H_2}\right). \tag{5.36}\label{eq:gr-multiplier}
\]
Choose an outer cutoff \(B_{\rm out}\), after the finite comparison
family has been fixed.
Dirichlet approximation with denominator ceiling
\(H_1H_2/L^{B_{\rm out}-1}\) makes the multiplier smaller than any
requested inverse logarithmic power unless \(T\theta\) lies within
\(L^{B_{\rm out}}/(H_1H_2)\) of a rational of denominator
\(r\le L^{B_{\rm out}}\).
The corresponding physical frequencies are covered by grids
of denominator \(q=Tr\), with widths
\(O(L^{B_{\rm out}}/(TH_1H_2))=O(L^{B_{\rm out}}/Y)\).
Moreover
\[
 \frac{Y}{qP_{\max}}
 \gg \exp(cL^{1/3}-CL^{1/4})L^{-B_{\rm out}},
\]
so \eqref{eq:gr-energy-ratio} holds uniformly in every gcd profile.

Discrete additive Parseval gives the exact normalization
\[
 \sum_{j\bmod q}|\widehat F(j/q+\beta)|^2
                         =qE_q(F\ee{\beta n}).        \tag{5.37}\label{eq:gr-gridparseval}
\]
Integrating over width \(L^{B_{\rm out}}/Y\) and applying
Proposition~\ref{prop:gr-energy} gives \(YL^{-D}\) for every desired
\(D\), after demanding a larger residue-energy saving.
The polynomial-logarithmic family of \(r\)'s is paid at the same step.
In particular the \(T\) grid points do not create a factor \(T\):
they are already included in \eqref{eq:gr-gridparseval}.
Cauchy--Schwarz and Parseval on the other endpoint bound
\eqref{eq:gr-fouriercomparison} on these grids by an arbitrary inverse
power of \(L\).  Off the grids, use
\eqref{eq:gr-multiplier} and the endpoint norms.
This proves that every comparison in \eqref{eq:gr-comparison} is
negligible to the requested logarithmic accuracy.

Finally \(\Psi(t)=t\Phi(t)\) makes
\[
               \frac{\Psi(|w|/x)}{|w|}
                       =\frac{\Phi(|w|/x)}x.
\]
Thus the harmonic raw correlation is \(1/x\) times the ordinary
banked correlation on the reflected line.
The comparison at scale \(Y\) has exactly the corresponding factor
\(1/Y\).  Paying the finite comparison coefficients therefore proves
the centered banked replacement required in the assembly.

Fixed bad-supported coefficients in \(Aa+Bb=T\) cause no new
analytic condition.  Embed, for example,
\(\widetilde F(n)=\1_{A\mid n}F(n/A)\).
Its exact residue energy is
\[
                  E_q(\widetilde F)=E_{q/(q,A)}(F),
\]
and its core length is divided by the fixed \(A\).
All energy ratios change only by fixed factors.
Bank primes are units at \(A,B\), so dilation invariance and
shared-label removal persist.  This embedding also avoids treating
the physical summands as units modulo \(L_0\) when they are not.

\section{Quantitative replacement and completion of the proof}\label{sec:assembly}
We now choose all accuracies in a fixed order. This is necessary
both for a positive logarithmic saving and to avoid feeding a later
density depth into the complexity of the earlier ideal construction.
Write
\[
 \phi=\varphi(L_0),\quad \delta=\frac1{8192\phi},\quad
 \mathcal S_N=\frac{NJ_{\mathcal E}(N)}{AB L^\kappa},\quad L=\log N.
\]
If \(\Gamma_N=0\), the count is zero by its fixed residue condition.
Hence assume \(\Gamma_N>0\); its finitely many possible positive
values have a fixed positive minimum.

\subsection{An explicit discarded-bank exponent}
The identity-class primes have reciprocal mass
\((1/(8\phi)+o(1))\log L\) in the small range and
\((1/(15\phi)+o(1))\log L\) in the big range. Greedily choose
\[
 b_N=\left\lfloor\frac{\log L}{32\phi}\right\rfloor
\]
disjoint banks of each size, each with reciprocal mass \(1+o(1)\).
There is ample mass in both ranges. The overshoot per bank is at
most the largest reciprocal prime, and target-prime deletion is
superlogarithmically small. Thus all masses are \(1+o(1)\), with
total error smaller than every inverse logarithmic power.
By \eqref{eq:success}, each success probability tends to
\(e^{-2}>1/8\). Chernoff's inequality and independence give
\begin{equation}\label{eq:rare-budget}
 \Prob\{\text{fewer than }b_N/16\text{ successes in either range}\}
 \le2e^{-b_N/64}\ll L^{-1/(4096\phi)}=L^{-2\delta}.
\end{equation}
Here \(b_N\ge\log L/(64\phi)\) eventually.
Lemma~\ref{lem:envelope} bounds the discarded correlation by
\(O(\mathcal S_NL^{-2\delta})\), plus exponential and power-smaller
errors. This applies with an indicator or model on either side.
The exact reconstruction \eqref{eq:reconstruction} has total
coefficient mass at most
\[
 3^{2b_N}\prod_g\max(1,V_g^2)\le L
\]
for large \(N\), since the second product is \(1+o(1)\) and
\(2b_N\log3\le(\log3)\log L/(16\phi)\).

\subsection{Moving cutoffs and the reflected target scale}
Set \(\epsilon=L^{-\delta}\). Choose \(0\le\Phi_\epsilon\le1\),
equal to one on \([2\epsilon,1-2\epsilon]\), supported in
\([\epsilon,1-\epsilon]\), with transitions of width \(\epsilon\).
For the harmonic graph comparison put
\(\Psi_\epsilon(t)=t\Phi_\epsilon(t)\) and
\(F_\epsilon(u)=\Psi_\epsilon(e^u)\). Direct differentiation gives
\[
 \|F_\epsilon\|_1\ll1,\qquad
 \|F_\epsilon^{(j)}\|_1\ll_j\epsilon^{1-j}\quad(j\ge1).
\]
At the lower transition the logarithmic derivative cancels the
small transition scale. Fourier integration by parts therefore gives
\begin{equation}\label{eq:mellin-budget}
 \int|\widehat F_\epsilon(\zeta)|(1+|\zeta|)^r\,d\zeta
       \ll_r\epsilon^{-r-1},\qquad
 \int_{|\zeta|>L^3}|\widehat F_\epsilon(\zeta)|\,d\zeta
       \ll_j L^{-(3-\delta)(j-1)}.
\end{equation}
Thus the external phase range stays fixed at \(L^3\); choose the
tail derivative order last.

On a contributing edge of the reflected graph,
\(n<0\), \(n'=n+ND>0\), and
\(t=|n|/(ND)\in[\epsilon,1-\epsilon]\). Exactly
\begin{equation}\label{eq:two-sided-localization}
 \frac{n'}{|n|}=\frac{1-t}{t}
       \in\left[\frac\epsilon{1-\epsilon},
                  \frac{1-\epsilon}{\epsilon}\right].
\end{equation}
For a source dyad \(|n|\asymp X\), insert a smooth target cutoff
supported on \([c\epsilon X,C\epsilon^{-1}X]\) and equal to one
on the contributing range \emph{before} Mellin inversion. Keep it
in every resulting pairing. Removed edges had zero full Mellin
integral originally. Dividing shared labels preserves the ratio
\eqref{eq:two-sided-localization}.

There are \(O(\log L)\) target dyads \(X'\), each with uniformly
bounded rescaled logarithmic derivatives. Their endpoint moments are
\(O(X)\), \(O(X')\): fixing an exact label product and counting its
multiples proves this directly, and summing reciprocal list mass
introduces no exponent depending on the final fixed mark length.
The pairing can therefore cost at most
\(\sqrt{X'/X}\ll\epsilon^{-1/2}\). Together with the unweighted
Mellin norm and dyadic count these costs are at most \(L^{5\delta/2}\).
We reserve \(L^{8\delta}\) for all localization losses.

For the comparison source scale \(Y\asymp ND_1D_2\), use its
actual target dyad \(Y'\), where
\(c\epsilon Y\le Y'\le C\epsilon^{-1}Y\). The physical major
window is \(L^{B_{\rm out}}/Y\), so
\(|\beta|Y'\ll L^{B_{\rm out}+\delta}\). Apply
Proposition~\ref{prop:gr-energy} with that enlarged exponent.
A crude integrated-energy rescaling costs \(\epsilon^{-2}\), still
inside the reserve. At the smallest dyad and selected small prime
\(P\le\exp(CL^{1/4})\),
\begin{equation}\label{eq:smallest-dyad}
 \frac{Y'}{qP}\gg
 e^{cL^{1/3}-CL^{1/4}}L^{-B_{\rm out}-\delta}\gg_A L^A.
\end{equation}
This remains true after every gcd split, since both length and
modulus are divided by the same divisor. Thus neither target
localization nor target valuations spoil the cofactor-energy estimate.

\subsection{The precision order}
Here are sufficient numerical margins. All exponents are fixed
before \(N\to\infty\).
First request normalized graph-pairing error \(L^{-22}\) per source
dyad. At most \(L^2\) source dyads occur, leaving \(L^{-20}\).
Fix the root and ideal accuracies for this request, then the probe
complexity and eventual mark length. The phase range is \(L^3\).

After that finite comparison family is fixed, let \(C_{\rm cmp}\)
dominate its coefficient and dyadic costs, and \(C_w\) its endpoint
pointwise exponent. These constants do not depend on the density
depth. Let \(m_*\) bound free draws per group. Unique factorization
gives a dyadic free-product law
\[
 \nu_i(n)\le L^{K_*}/n,\qquad
 K_*=\frac{\log(m_* !)+m_*\log2}{16\phi}.
\]
Since the law has mass at most one, its squared norm on a dyad of
length \(T_i\) is at most \(L^{K_*}/T_i\). The integer bilinear
bound \eqref{eq:gr-multiplier} consequently permits
\(C_{\rm prod}=2K_*+2\). Choose
\begin{align}
 A_{\min}&>C_{\rm cmp}+2C_w+40+8\delta,\notag\\
 B_{\rm out}&>2A_{\min}+C_{\rm prod}+4.\label{eq:outer-budget}
\end{align}
Dirichlet approximation with denominator ceiling
\(T_1T_2/L^{B_{\rm out}-1}\) yields a multiplier
\(O(L^{-A_{\min}})\) off the grids of denominators \(Nr\),
\(r\le L^{B_{\rm out}}\). The middle terms in
\eqref{eq:gr-multiplier} are superlogarithmically small because
\(T_i\ge\exp(L^{1/3})\).

Next request integrated target-grid energy saving
\[
 T_E>2(C_{\rm cmp}+C_w+40+8\delta).
\]
Use window exponent \(B_{\rm out}+\delta\), and choose its
residue-energy accuracy to pay both the window and the number of
\(r\)'s. Parseval on the other endpoint then gives major contribution
at most \(L^{-T_E/2+C_w+C_{\rm cmp}+8\delta}\le L^{-40}\).
The minor contribution has the same margin by
\eqref{eq:outer-budget}. Exact discrete Parseval, not a supremum
bound, cancels the target-grid factor \(N\).

Finally choose the Hankel depth for the required exceptional-character
centering. If \(C_{\rm low}\) pays record, character and cutoff costs,
take \(3(K+1)/5>T_E+C_{\rm low}+40\).
Choose the derivative order in \eqref{eq:mellin-budget} after its
tail costs. These choices change constants and thresholds, not the
earlier ideal complexity, bounded endpoint moment exponent or phase
range. In particular the precision choices are noncircular.

\begin{proof}[Completion of Theorem~\ref{thm:general}]
After localization, the graph and comparison errors are bounded by
\(O(NL^{-20+8\delta})\). Reconstruction costs at most \(L\), so
the resulting error is
\[
 O(NL^{-19+8\delta})=o(\mathcal S_NL^{-\delta}),
\]
using \(\kappa<2\), \(J_{\mathcal E}\ge1\), and fixed \(A,B\).
The discarded-bank error is \(O(\mathcal S_NL^{-2\delta})\).
Apply the resulting replacement on each side to the identity
\[
 f_U(x)f_V(y)-D_U(x)D_V(y)
   =\xi_U(x)f_V(y)+D_U(x)\xi_V(y).
\]
Proposition~\ref{prop:model-correlation} evaluates the remaining
model correlation. On retained bulk records its logarithm is
\(L+O(\sqrt L+\delta\log L)\), giving relative
\(O(L^{-1/2})\) in leading-density replacement. The sieve uses a
fixed power level; partial summation costs
\(\TV(\Phi_\epsilon)=O(1)\), rather than its largest derivative.
The record and squared-modulus CRT endpoints cost
\(N^{0.45+o(1)}\). The upper endpoint sieve gives
\[
 O(\epsilon\mathcal S_N+N^{0.45+o(1)}).
\]
Its strip length \(N/L^\delta\) exceeds the sieve level. Euler tails
are superlogarithmically small, and
\(\int\Phi_\epsilon=1+O(\epsilon)\). Every error is therefore
\(O(\mathcal S_NL^{-\delta})\), as asserted.

For an interval \(I\), if \(|I|\ll\epsilon\) use the direct upper
sieve. Otherwise clip it to \([\epsilon,1-\epsilon]\) and smooth
its endpoints at width \(\epsilon\). At most four transition strips
of total length \(O(\epsilon N)\) are lost. Total variation remains
\(O(1)\), and the logarithmic derivative norm of a transition at
\(t_0\) is at most \(O_j(t_0^j\epsilon^{1-j})\).
Thus the same budgets apply uniformly in both endpoints of \(I\),
and the cutoff integral is \(|I|+O(\epsilon)\).
Endpoint conventions alter at most two pairs. This proves the
uniform interval assertion and completes the theorem.
\end{proof}

\subsection{Principal norm specialization}\label{sec:norms}

We recover primitive principal norm values by taking both forms in
the preceding results to be the norm form of one order. Thus the
ideal-class comparison and the local-density normalization require no
separate argument.

Fix an imaginary quadratic field $K$ of fundamental discriminant
$D<0$, and an order $\mathcal O=\Z+f\mathcal O_K$ of discriminant
$\Delta=f^2D$.  Choose a basis $1,\omega$ of $\mathcal O$ and put
\[
 Q(x,y)=\operatorname N_{K/\Q}(x+y\omega),\qquad
 P=\{Q(x,y)>0:x,y\in\Z,\ \gcd(x,y)=1\}.
\]
Primitivity refers to this order basis.  Write
\[
 \begin{gathered}
 \mathscr V_p=Q(\Z_p^2\setminus p\Z_p^2),\qquad
 \Lambda=\{n\ge1:n\in\mathscr V_p\text{ for every }p\},\\
 R_T(N)=\#\{(a,b)\in T^2:a+b=N\}.
 \end{gathered}
\]
Thus $P\subset\Lambda$, and $R_T$ counts numerical values without
coordinate multiplicities.  A target is locally admissible if
$N\in\mathscr V_p+\mathscr V_p$ for every prime $p$.

Let $S=\{p:p\mid2\Delta\}$, let $\chi=\chi_D$, and define
\[
 d_p(N)=\int_{\Z_p}\mathbf1_{\mathscr V_p}(a)
                    \mathbf1_{\mathscr V_p}(N-a)\,da\quad(p\in S),
 \qquad
 J_{\mathcal O}(N)=
 \prod_{\substack{p\mid N\\p\notin S,\ \chi(p)=-1}}
                         \frac{p-1}{p-2}.
\]
The Haar measure of $\Z_p$ is one.  The local densities concern the
set of values, rather than the number of their preimages under $Q$.

An independent weighted quadratic correlation and positive restoration
give an order-of-magnitude result. Its classical core estimate is proved
in full in Appendix~\ref{sec:older-route}, using no new zero-free or
weighted-dilation input.

\begin{corollary}[The quadratic order bounds]\label{thm:norm-order}
For every fixed imaginary quadratic order there are constants
$c_{\mathcal O},C_{\mathcal O}>0$ and $N_{\mathcal O}$ such that every
locally admissible $N\ge N_{\mathcal O}$ satisfies
\[
 c_{\mathcal O}\frac{NJ_{\mathcal O}(N)}{\log N}
       \le R_P(N)\le R_\Lambda(N)
       \le C_{\mathcal O}\frac{NJ_{\mathcal O}(N)}{\log N}.
\]
Every nonadmissible target has count zero.  These bounds use
Lemma~\ref{lem:older-quadratic-filter} and
Proposition~\ref{prop:norm-class-error}; they do not use
Theorem~\ref{thm:general}.
\end{corollary}

The leading-term theorem gives the following explicit specialization.

\begin{corollary}\label{thm:norm-leading}
Under the analytic inputs of Section~\ref{sec:inputs}, for every fixed
imaginary quadratic order, uniformly as $N$ tends to infinity through
the locally admissible targets,
\[
 R_P(N)\sim\mathfrak A_{\mathcal O}(N)
                     J_{\mathcal O}(N)\frac N{\log N},
\]
where the explicit positive coefficient is
\begin{equation}\label{eq:norm-coefficient}
 \mathfrak A_{\mathcal O}(N)=
 \frac{4L(1,\chi)}{\pi}
 \prod_{p\in S}\frac{1-\chi(p)/p}{1-1/p}
 \prod_{\substack{p\notin S\\\chi(p)=-1}}
                  \left(1-\frac2{p(p-1)}\right)
 \prod_{p\in S}d_p(N).
\end{equation}
Every nonadmissible target has count zero.  Each $d_p(N)$ is a
computable rational periodic function, and the infinite product
converges absolutely.  No restriction on the conductor or the ring
class group is required; the order remains fixed.
\end{corollary}

\begin{proof}[Specialization]
The norm form is primitive of discriminant \(\Delta\) and corresponds
to \(A=\mathcal O\), hence \(c_Q=0\). Lemma~\ref{lem:form-ideal-shadow}
therefore reads
\begin{equation}\label{eq:norm-shadow-ideals}
 n\in\Lambda\quad\Longleftrightarrow\quad
 \text{some primitive integral invertible }I
       \text{ has }n(I)=n,\ [I]\in2G,
\end{equation}
and replacing \(2G\) by zero characterizes \(P\).
The classical bounds just proved for mixed forms give
Corollary~\ref{thm:norm-order}. The two local filters now coincide,
\(d_{p;1,2}=d_p\), and \(J_\Delta=J_{\mathcal O}\), so
Theorem~\ref{thm:form-binary-leading} gives
\eqref{eq:norm-coefficient} and Corollary~\ref{thm:norm-leading}.
In particular, Proposition~\ref{prop:norm-class-error} gives
\(R_\Lambda-R_P\ll_{\mathcal O}NJ_{\mathcal O}/(\log N)^{6/5}\);
if \(2G=0\), it gives \(P=\Lambda\) exactly.
The finite local computation, prefix upper bound, summability and
Haar normalization are given by
\eqref{eq:norm-local-density-finite},
\eqref{eq:norm-prefix-tail} and \eqref{eq:form-haar-strata}.
They include all conductor-local factors in the principal case.
\end{proof}

\begin{remark}[Gaussian normalization]
For $\mathcal O=\Z[i]$ one has $P=C_4\sqcup2C_4$ and
\[
 d_2(N)=
 \begin{cases}
 1/4,&N\bmod8\in\{2,3,6,7\},\\
 1/8,&N\bmod8=4,\\
 0,&N\bmod8\in\{0,1,5\}.
 \end{cases}
\]
Consequently the coefficient of $J_{\mathcal O}(N)N/\log N$ is
$\frac12\prod_{p\equiv3\ (4)}(1-2/(p(p-1)))$ in the first four
classes and half that coefficient in the class $4\bmod8$.
This also calibrates the value-density normalization: the two Gamma
factors are already included in $c_{\mathcal C}^2$, and no additional
beta-integral factor occurs.
\end{remark}

\section{Consequences and boundaries}
\begin{corollary}\label{cor:basis13}
The set \(C_{12}\cup\{0\}\) is an asymptotic basis of exact order
thirteen: every sufficiently large integer is a sum of at most
thirteen elements of \(C_{12}\), and twelve do not suffice.
\end{corollary}
\begin{proof}
For a large target choose \(j\in\{0,\ldots,11\}\) with
\(N-j\equiv2\pmod{12}\). Theorem~\ref{thm:c12} writes \(N-j\)
as two core values; add \(j\) copies of one. This uses at most thirteen
positive summands. Conversely, \(25+60r\), \(r\ge0\), is
\(1\pmod{12}\) and divisible by the forbidden prime five.
At most twelve positive core summands would force exactly one
summand on that residue class, which is impossible. Padding with
zeros gives the equivalent fixed-order assertion.
\end{proof}

The field formulation follows from classical abelian class field
theory. A subgroup condition on Frobenius is complete splitting in
its fixed field. Finite conjunctions become intersections of inverse
image subgroups in the compositum. Kronecker--Weber places this
compositum in a cyclotomic field
\cite[Chapter I, Theorem 4.16]{milnecft}, giving precisely the
fixed subgroups in Theorem~\ref{thm:general}, after discarding the
fixed bad-prime set. The reduction is classical; the pointwise
indicator correlation is the analytic assertion used here.

These conclusions have several precise boundaries. The semigroup
theorems treat subgroup pairs and unfiltered self-sums for a prime
coset. Other unions of residue classes require their own model and
sieve checks. The modulus and coefficients stay fixed. The interval discrepancy is absolute,
and gives no relative asymptotic in intervals of power-short length
\(N^\theta\), \(\theta<1\). The threshold is not numerical, so it
does not classify small exceptions. The exact sumsets of primitive
quadratic norms and their finite certificates form a separate
arithmetic problem, treated in the companion article
\emph{Exact Sumsets of Primitive Quadratic Norms}.
The separate article
\href{https://hexagonmath.org/2611.00593v1}{\emph{Pointwise representation
counts for Eisenstein restricted sums}} addresses where summands may
be placed, with a different proof and parameter range.

The shared comparison requires a common invariant prime class.
Two arbitrary prime prescriptions can be disjoint, so that no such
class is available. Their binary correlation is not a consequence of
Theorem~\ref{thm:model-comparison}. Nor does its fixed-shift clause
assert uniformity for growing sublinear shifts between its two
proved frequency regimes.

Higher-fold sums and independent coupled systems can use classical
linear-form correlation theorems, but their proof structure differs
from this binary reflected line. The ordinary quadratic-form
indicator correlations in Section~\ref{sec:quadratic-pairs} use
the fixed-shift clause after separately verifying its arithmetic
hypotheses. They include the Connors--Keating case and Smilansky's
five-form family. Growing-shift uniformity and higher indicator
correlations remain outside those statements.

\paragraph{AI assistance and review.}
ChatGPT (OpenAI) assisted with mathematical arguments, literature checks,
exposition, proof review and verification code.
\begin{thebibliography}{99}\small

\bibitem{bourgainfuchs}
J. Bourgain and E. Fuchs,
\emph{On representation of integers by binary quadratic forms},
arXiv:1105.3698v2, 23 May 2011.
\url{https://arxiv.org/abs/1105.3698v2}.

\bibitem{cox}
D. A. Cox,
\emph{Primes of the Form \(x^2+ny^2\): Fermat, Class Field Theory,
and Complex Multiplication}, second edition, Wiley, 2013.
\url{https://doi.org/10.1002/9781118400722}.

\bibitem{delabretechetenenbaum}
R. de la Bret\`eche and G. Tenenbaum,
\emph{Remarks on the Selberg--Delange method},
Acta Arithmetica \textbf{200} (2021), 349--369.
\url{https://doi.org/10.4064/aa201024-26-5}.
Author's corrected text:
\url{https://tenenb.perso.math.cnrs.fr/PPP/On-SD.pdf}.

\bibitem{halberstamrichert}
H. Halberstam and H.-E. Richert,
\emph{Sieve Methods}, London Mathematical Society Monographs 4,
Academic Press, London--New York, 1974.

\bibitem{iwaniecpomykala}
H. Iwaniec and J. Pomyka{\l}a,
\emph{Sums and differences of quartic norms},
Mathematika \textbf{40} (1993), 233--245.
\url{https://doi.org/10.1112/S0025579300007014}.

\bibitem{khayutin}
I. Khayutin,
\emph{Joint equidistribution of CM points},
Annals of Mathematics (2) \textbf{189} (2019), 145--276.
\url{https://doi.org/10.4007/annals.2019.189.1.4}.

\bibitem{loughranmatthiesen}
D. Loughran and L. Matthiesen,
\emph{Frobenian multiplicative functions and rational points in fibrations},
Journal of the European Mathematical Society \textbf{26} (2024),
4779--4830.
\url{https://doi.org/10.4171/JEMS/1374}.

\bibitem{matthiesen}
L. Matthiesen,
\emph{Linear correlations of multiplicative functions},
Proceedings of the London Mathematical Society (3)
\textbf{121} (2020), 372--425.
\url{https://doi.org/10.1112/plms.12309}.
Theorem numbering here follows arXiv:1606.04482v4:
\url{https://arxiv.org/abs/1606.04482v4}.

\bibitem{milnecft}
J. S. Milne,
\emph{Class Field Theory}, version 4.03, 6 August 2020.
\url{https://www.jmilne.org/math/CourseNotes/CFT.pdf}.

\bibitem{weightedgraphs}
OpenAI,
\emph{Weighted dilation graphs, smooth shifted primes and totient fibers},
24 September 2026, Section 3.
\url{https://github.com/openai/math/tree/main/preprints/Weighted-Dilation-Graphs-Smooth-Shifted-Primes-and-Totient-Fibers-September-24-2026}.

\bibitem{zerofree}
OpenAI,
\emph{The Quasi-Riemann Hypothesis:
A Zero-Free Half-Plane \(\Re(s)>7/8\)},
30 September 2026, Part I, Theorem 3.1.
\url{https://github.com/openai/math/tree/main/preprints/The-Quasi-Riemann-Hypothesis-September-30-2026}.

\bibitem{ramareviswanadham}
O. Ramar\'e and G. K. Viswanadham,
\emph{The trigonometric polynomial on sums of two squares,
an additive problem and generalisation},
arXiv:2509.23260v1, 27 September 2025;
Transactions of the American Mathematical Society.
\url{https://doi.org/10.1090/tran/9598}.
\url{https://arxiv.org/abs/2509.23260v1}.

\bibitem{scharlau}
W. Scharlau,
\emph{Quadratic and Hermitian Forms},
Grundlehren der mathematischen Wissenschaften 270,
Springer, Berlin--Heidelberg, 1985.
\url{https://doi.org/10.1007/978-3-642-69971-9}.

\bibitem{serre}
J.-P. Serre,
\emph{A Course in Arithmetic},
Graduate Texts in Mathematics 7, Springer, New York, 1973.
\url{https://doi.org/10.1007/978-1-4684-9884-4}.

\bibitem{walsh}
M. N. Walsh,
\emph{Shifted convolutions of divisor functions},
arXiv:2610.40439v2, 12 November 2026.
\url{https://arxiv.org/abs/2610.40439v2}.

\end{thebibliography}

\end{document}
